\documentclass[11pt]{article}
\usepackage[utf8]{inputenc}
\usepackage[T1]{fontenc}
\usepackage{lmodern}
\usepackage{geometry}
\geometry{margin=1.2in}
\usepackage{amsmath, amsthm, amssymb, mathtools}
\usepackage{enumitem}
\usepackage{tikz-cd}
\usepackage{quiver}
\usepackage[colorlinks=true, linkcolor=blue, citecolor=red, urlcolor=blue]{hyperref}
% \usepackage[style=numeric, backend=biber, giveninits=true]{biblatex}
% \addbibresource{references.bib}
\bibliographystyle{abbrvurl}
% \renewcommand{\bibfont}{\footnotesize}
\usepackage{xurl}

\author{David Forsman\\
Université catholique de Louvain\\
\texttt{david.forsman@uclouvain.be}}
\title{A Syntactic Approach to Ulmer's Bialgebras}
\date{21 July, 2026}

\theoremstyle{plain}
\newtheorem{theorem}{Theorem}[section]
\newtheorem{lemma}[theorem]{Lemma}
\newtheorem{corollary}[theorem]{Corollary}
\newtheorem{proposition}[theorem]{Proposition}


\theoremstyle{definition}
\newtheorem{definition}[theorem]{Definition}
\newtheorem{example}[theorem]{Example}
\newtheorem{examples}[theorem]{Examples}
\newtheorem{remark}[theorem]{Remark}

\numberwithin{equation}{section}


\newcommand{\Mod}{\mathbf{Mod}}
\newcommand{\Alg}{\mathbf{Alg}}
\newcommand{\C}{\mathcal{C}}
\newcommand{\D}{\mathcal{D}}
\newcommand{\E}{\mathcal{E}}
\newcommand{\K}{\mathcal{K}}
\newcommand{\Term}{\mathrm{Term}}
\newcommand{\Cont}{\mathrm{Cont}}
\newcommand{\Obj}{\mathrm{Obj}}
\newcommand{\Mor}{\mathrm{Mor}}
\newcommand{\Const}{\mathrm{Const}}
\newcommand{\Hopf}{\mathbf{Hopf}}
\newcommand{\Set}{\mathbf{Set}}
\newcommand{\Cat}{\mathbf{Cat}}
\newcommand{\CAT}{\mathbf{CAT}}
\newcommand{\ACC}{\mathbf{ACC}}
\newcommand{\Ins}{\mathrm{Ins}}
\newcommand{\id}{\mathrm{id}}
\newcommand{\Eq}{\mathrm{Eq}}
\newcommand{\cod}{\mathrm{com}}
\newcommand{\dom}{\mathrm{dom}}
\newcommand{\LE}{\mathrm{LE}}
\newcommand{\RE}{\mathrm{RE}}
\newcommand{\V}{\mathcal{V}}
\newcommand{\Mon}{\mathbf{Mon}}
\newcommand{\colim}{\mathrm{colim}}

\begin{document}

\maketitle

\begin{abstract}
Ulmer introduced a semantic notion of bialgebras that unifies a broad class of algebraic and coalgebraic structures. We develop a syntactic counterpart by introducing signature pairs $(\Sigma,\sigma)$ and bialgebraic theories $T$, providing a uniform language for constructing internal bialgebras in a $2$-categorical setting. For every bialgebraic theory $T$ and $\Sigma$-model $M$ within a $2$-category with PIE limits, we construct the object $M^T$ of internal $T$-bialgebras. Our approach to bialgebras admits a general Induced Functor of Algebras Theorem extending the classical lifting of lax monoidal functors to the categories of internal monoids. Since the construction of $M^T$ is expressed entirely in terms of PIE limits, accessibility, local presentability, orthogonal factorization systems, regularity, and exactness lift along the construction $M \mapsto M^T$ under suitable assumptions.
\end{abstract}
\section{Introduction}

In his 1977 preprint, Ulmer \cite{Ulmer1977Bialgebras} introduced a semantic notion of a bialgebra to unify a broad range of algebraic and coalgebraic structures within the setting of locally presentable categories. The framework encompasses examples ranging from classical structures in universal algebra, such as groups and rings, and Eilenberg-Moore categories of monads, to bialgebraic structures arising in monoidal categories, such as internal bimonoids and Hopf algebras. Rather than treating these examples individually, Ulmer showed that many of their fundamental structural properties follow from a common semantic framework.

The present paper develops a syntactic counterpart to Ulmer's theory of bialgebras. The guiding idea is that the support functors (domains and codomains of structural maps) and natural transformations (structural equations) appearing in Ulmer's definition of a bialgebra should arise automatically from logical syntax rather than being specified individually for each notion of a bialgebra. To achieve this, we introduce multi-sorted signature pairs $(\Sigma,\sigma)$ with associated bialgebraic theories $T$. The $2$-signature $\Sigma$ fixes the structure of a $\Sigma$-model $M$, for instance, a monoidal or monadic structure, while the relative multi-sorted $1$-signature $\sigma$ and the bialgebraic theory $T$ describe the operations and equations defining the structure $M^T$ of the bialgebras within $M$. This produces a syntactic language whose interpretation in categories produces exactly the structures of bialgebras introduced by Ulmer, while automating the construction of the associated support functors and natural transformations from the syntax itself.

The central observation of this paper is that the commutative diagrams defining bialgebraic structures can be interpreted both as abstract syntax and as objectwise equations. Thus, the same diagrams that define, for example, Frobenius algebras in a monoidal category, determine a bialgebraic theory which admits an interpretation within any pseudo monoid in a $2$-category with PIE limits (products, inserters, and equifiers). We show that, for every bialgebraic theory $T$, the corresponding object $M^T$ of internal $T$-bialgebras is obtained entirely by means of strict PIE limits. Consequently, we generalize the construction $M\mapsto M^T$ from the $2$-category $\CAT$ to arbitrary $2$-categories with PIE limits. What is perhaps surprising is that the theory $T$ consists of parallel paths, representing commutative diagrams, and yet we are able to assign meaning to the notion of $T$-bialgebras even in $2$-categories whose $0$-cells do not have a clear notion of objects nor of morphisms. This framework also yields a general induced-functor theorem extending the classical lifting of lax monoidal functors to categories of internal monoids.

Expressing the construction of $M^T$ entirely in terms of PIE limits has several important structural consequences. A recent result by Lack and Tendas \cite[Theorem~5.5.]{LACK2023107196} shows that the $2$-category $\V$-$\ACC$ of $\V$-accessible categories, $\V$-accessible functors, and $\V$-natural transformations is closed under the flexible $2$-limits in $\V$-$\CAT$ of $\V$-categories. As PIE limits are flexible \cite{PowerPIE1991}, this immediately implies that, whenever the $\Sigma$-model $M$ consists of accessible categories and accessible functors, the object $M^T$ is again accessible, both in the ordinary and enriched settings.\footnote{The base $\V$ of enrichment is assumed to be a locally presentable symmetric monoidal closed category.} This extends recent accessibility results for particular classes of bialgebras \cite{agore2025categoryhopfbraces, PorstHopfLocallyPresentable, Porst01062008}.

To obtain more structural properties for the object $M^T$ of bialgebras, we prove that the canonical morphisms arising from PIE limits strictly create left and right Kan extensions under suitable assumptions. With some care, this allows local presentability to lift along the construction
$$
M \longmapsto M^T.
$$
Independently, we establish conditions under which orthogonal factorization systems lift along PIE limits, from which we derive corresponding lifting results for regularity and exactness.

The paper is organized as follows. Section~2 reviews the required $2$-categorical preliminaries. We define strict $2$-categories emphasizing strict PIE limits, which serve as the fundamental construction throughout the paper. Section~3 develops the syntax and semantics of signature pairs, culminating in the Induced Functor of Algebras Theorem (Theorem~\ref{thm:induced_morphism_between_algebras}), which establishes conditions under which a lax, oplax, or strong morphism $M\to N$ of $\Sigma$-models induces a morphism $M^T\to N^T$ between the corresponding objects of bialgebras. Section~4 establishes conditions for the strict creation of left and right Kan extensions along PIE limits (Theorem~\ref{thm:pie_kan_creation}), applies the recent work of Lack and Tendas to prove the lifting of accessibility (Theorem~\ref{thm:enriched_accessability_lifting}), and combines these results to obtain the lifting of local presentability (Theorem~\ref{thm:LiftingLP}). Finally, Section~5 studies the lifting of orthogonal factorization systems along PIE limits (Theorem~\ref{thm:pie_ofs_creation}) and applies this result to derive corresponding lifting theorems for regularity and exactness, culminating in Corollary~\ref{cor:StructuresLiftedToBialgebras}.

\section{2-Categories}
Since the construction of internal bialgebras developed in this paper is expressed entirely in terms of strict PIE limits, we begin by establishing the required $2$-categorical framework. We briefly review enriched categories, concentrating on the case where the enriching category is the category of locally small categories, thereby recovering the standard notion of a $2$-category. For a comprehensive treatment of enriched category theory, we refer the reader to \cite{Kelly:Enriched}.

We then recall the definitions of strict products, inserters, and equifiers, collectively known as the strict PIE limits. These $2$-limits form the basic building blocks for all subsequent constructions. Throughout the paper, we work in a strict $2$-categorical setting. This entails no essential loss of generality, since strict PIE limits are flexible limits, and therefore their constructions transport to bicategorical settings via the bicategorical strictification \cite{PowerPIE1991, GordonPowerStreet1995CoherenceTricategories}. 

\subsection{Enriched Categories and 2-Categories}

\begin{definition}[$\V$-Category]
   Let $\V = (V,\otimes, I)$ be a monoidal category. A $\V$-enriched category $\K$ consists of the following data:\footnote{We assume the existence of a Grothendieck universe $U$ and its successor universe $U^+$. The elements and subsets of $U$ are called small and large, respectively, and the elements of $U^+$ are called very large. Note that small sets are both large and very large.}
   \begin{itemize}
       \item A set $S$ of \textbf{objects}, often denoted $a,b,c,\dots$.
       \item For any pair of objects $a, b\in S$, a \textbf{hom-object} $\K(a,b)$ in $\V$.
       \item For any objects $a, b, c$, a \textbf{composition morphism} $* = *_{a,b,c}\colon \K(b,c) \otimes \K(a,b) \to \K(a,c)$ in $\V$ and a chosen identity $\id_a\colon I\to \K(a,a)$. Furthermore, we require that the composition is associative and $\id_a$ defines the identity for each $a\in S$.
   \end{itemize}
   When $\V$ is the cartesian closed category of very large categories, a very large $\V$-category $\K$ is called a $2$-\textbf{category}.\footnote{We call a $\V$-category small/large/very large if the underlying set of objects is.} We denote the $2$-category of large $\V$-categories by $\V$-$\CAT$. Unpacking this $\V$-enriched definition, a $2$-category $\K$ is equipped with three forms of composition:
   \begin{enumerate}
       \item \textbf{Vertical composition} of transformations: A morphism in $\K(a,b)$ is denoted $\eta\colon f\Rightarrow g\colon a\to b$, where $f,g$ and $\eta$ are called morphisms and transformations of $\K$, respectively. For transformations $\alpha \colon f \Rightarrow g\colon a\to b$ and $\beta \colon g \Rightarrow h\colon a\to b$ in $\K$, their composite in $\K(a,b)$ is written $\beta \bullet \alpha \colon f \Rightarrow h$.
       \item \textbf{Horizontal composition} of morphisms: For $f \colon a \to b$ and $g \colon b \to c$, their composite is denoted $gf \colon a \to c$.
       \item \textbf{Horizontal composition} of transformations: For $\alpha \colon f \Rightarrow g\colon a\to b$ and $\beta \colon h \Rightarrow k\colon b\to c$, their composite is denoted $\beta * \alpha \colon hf \Rightarrow kg\colon  a\to c$. Note that the bifunctoriality of $*$ yields the interchange equations:
       $$
       \beta*\alpha = \beta g\bullet h\alpha = k\alpha\bullet \beta f
       $$where we use the whiskering notation $\beta h\coloneqq \beta*\id_h$ and $h\alpha\coloneqq \id_h*\alpha$. We identify the functor $a\colon 1\to \K(a,a)$ with the morphism $a(!)$, where $!$ is the unique object of the terminal category $1$. 
   \end{enumerate}
\end{definition}

\subsection{PIE limits}
We define the so-called strict PIE limits (Products, Inserters, and Equifiers).
\begin{definition}[PIE Limits in a $2$-Category]
   Let $\K$ be a $2$-category. We define the following $2$-dimensional limits by their explicit universal properties on both $1$-cells and $2$-cells:
   \begin{itemize}
       \item \textbf{Products:} The product of a family of objects $(c_i)_{i \in I}$ is an object $\prod_i c_i$ equipped with projection morphisms $\pi_j \colon \prod_i c_i \to c_j$ satisfying:
       \begin{enumerate}[label=(\roman*)]
           \item For any family of morphisms $(f_i\colon x\to c_i)_{i\in I}$ in $\K$ there is a unique morphism $f\colon x\to \prod_i c_i$ so that $f_i = \pi_i\circ f$. We often denote $f = (f_i)_i$.
           \item For any family $(\alpha_i\colon f_i\Rightarrow g_i\colon x\to c_i)_{i\in I}$ of transformations in $\K$, there exists a unique transformation $\alpha\colon (f_i)_i\Rightarrow  (g_i)_i\colon x\to \prod_i c_i$ so that $\alpha_i = \pi_i* \alpha$, often denoted $\alpha = (\alpha_i)_i$. 
       \end{enumerate}
       
       \item \textbf{Inserters:} Let $f,g\colon c\to d$ be parallel morphisms in $\K$. The \textbf{inserter} of $f$ and $g$ is an object $\Ins(f,g)$ in $\K$ equipped with a morphism $u\colon\Ins(f,g) \to c$ and a universal transformation $\phi \colon f \circ u \Rightarrow g \circ u$ satisfying:
       \begin{enumerate}[label=(\roman*)]
           \item For any morphism $h'\colon x\to c$ with a transformation $\alpha\colon fh'\Rightarrow gh'$, there exists a unique morphism $h\colon x\to \Ins(f,g)$ so that $h' = uh$ and $\alpha = \phi*h$. 
           \item For any pair of morphisms $h_1,h_2\colon x\to \Ins(f,g)$ with a transformation $\alpha'\colon uh_1\Rightarrow uh_2$ making the diagram 
           $$
           \begin{tikzcd}
fuh_1 \arrow[d, "f*\alpha'"', Rightarrow] \arrow[r, "\phi*h_1", Rightarrow] & guh_1 \arrow[d, "g*\alpha'", Rightarrow] \\
fuh_2 \arrow[r, "\phi*h_2"', Rightarrow]                                   & guh_2                                  
\end{tikzcd}
           $$
           commute, there is a unique transformation $\alpha\colon h_1\Rightarrow h_2$ so that $\alpha' = u*\alpha$.
       \end{enumerate}
       
       \item \textbf{Equifiers:} Let $\eta, \theta \colon f\Rightarrow g\colon c\to d$ be parallel transformations in $\K$. The \textbf{equifier} of $\eta$ and $\theta$ is an object $\Eq(\eta, \theta)$ equipped with a morphism $u \colon \Eq(\eta, \theta) \to c$ such that $\eta * u = \theta * u$, satisfying:
       \begin{enumerate}[label=(\roman*)]
           \item Any morphism $h'\colon x\to c$, satisfying $\eta*h' = \theta*h'$, has a unique lift $h\colon x\to \Eq(\eta,\theta)$ along $u$, meaning $h' = uh$.
           \item The functor $u*(-)\colon \K(x,\Eq(\eta,\theta))\to \K(x,c)$ is fully faithful for every object $x$ in $\K$.
       \end{enumerate}
   \end{itemize}
   We say that $\K$ is \textbf{PIE-complete} if it has small products, inserters and equifiers.
\end{definition}

\begin{remark}[PIE Limits in $\CAT$]
   In the $2$-category $\CAT$, small products are standard cartesian products of categories, and inserters and equifiers are computed explicitly:
   \begin{itemize}
       \item Let $F, G \colon \C \to \mathcal{D}$ be parallel functors. The \textbf{inserter category} $\operatorname{Ins}(F, G)$ has as objects pairs $(c, \phi\colon Fc\to Gc)$, where $c$ is an object of $\C$ and $\phi$ is a morphism in $\mathcal{D}$. A morphism $f\colon (c,\alpha)\to (d,\beta)$ in $\Ins(F,G)$ consists of a morphism $f\colon c\to d$ in $\C$ making the diagram
       $$
       \begin{tikzcd}
Fc \arrow[d, "F(f)"'] \arrow[r, "\alpha"] & Gc \arrow[d, "G(f)"] \\
Fd \arrow[r, "\beta"']                             & Gd                  
       \end{tikzcd}
       $$
       commute. The universal functor $U \colon \Ins(F, G) \to \C$ is the forgetful functor $f\colon (c,\alpha)\to (d,\beta)\mapsto f\colon c\to d$, and the universal natural transformation is $\phi= (\alpha)_{(c,\alpha)\in\mathrm{Obj}(\Ins(F,G))}$.

       \item Let $\eta,\theta\colon F\Rightarrow G\colon \C\to \mathcal{D}$ be natural transformations. The \textbf{equifier category} $\operatorname{Eq}(\eta, \theta)$ is the full subcategory of $\C$ consisting of those objects $c \in \C$ for which the components agree $\eta_c = \theta_c$. The universal functor $E \colon \operatorname{Eq}(\eta, \theta) \to \C$ is the canonical fully faithful inclusion.
   \end{itemize}
\end{remark}
It is perhaps useful to notice that PIE limits are defined exactly so that representable $2$-functors $\K(x,-)\colon \K\to \CAT$ will respect PIE limits, where the PIE limits in $\CAT$ are a fixed construction prior to the notion of $2$-limits.


\section{Bialgebras}

In this section, we develop a multi-sorted, $2$-categorical syntactic framework to study internal algebraic and coalgebraic structures. Conceptually, our formulation of a bialgebra specialized in the $2$-category of $\CAT$ is equivalent to the notion introduced by Ulmer \cite{Ulmer1977Bialgebras}, where algebraic structures are captured by inserters and equations are enforced by equifiers. 

Specifically, for a category $\underline{A}$ equipped with an $M$-indexed family of functors pairs
$$
(F_{d\mu},F_{c\mu}\colon \underline{A}\to X_\mu)_{\mu\in M},
$$
Ulmer defined the category of pre-bialgebras $\mathrm{P\text{-}Bialg}_M(\underline{A})$ as the inserter $\Ins(F_d,F_c)$ of the product induced functors $F_d\coloneqq (F_{d\mu})_\mu$ and $F_c\coloneqq (F_{c\mu})_\mu \colon \underline{A}\to \prod_{\mu\in M} X_\mu$. Denoting the canonical forgetful functor by $V\colon \mathrm{P\text{-}Bialg}_M(\underline{A})\to \underline{A}$, and given an $R$-indexed family of parallel pairs of natural transformations:
$$
\begin{tikzcd}
\mathrm{P\text{-}Bialg}_M(\underline{A}) \arrow[d, "V"'] \arrow[r, "V"]                                                           & \underline{A} \arrow[d, "F_{cr}"] \\
\underline{A} \arrow[r, "F_{dr}"'] \arrow[ru, "\alpha_r", Rightarrow, shift left] \arrow[ru, "\beta_r"', Rightarrow, shift right] & X_r                              
\end{tikzcd}
$$ 
the category $\mathrm{Bialg}_{M,R}(\underline{A})$ of bialgebras is defined as the equifier $\Eq(\alpha,\beta)$ of the parallel natural transformations $\alpha\coloneqq (\alpha_r)_r$ and $\beta\coloneqq (\beta_r)_r\colon (F_{dr} V)_r\Rightarrow (F_{cr}V)_r\colon \mathrm{P\text{-}Bialg}_M(\underline{A})\to \prod_{r\in R}X_r$.

While this formulation is purely $2$-categorical, it requires manually constructing the support functors and natural transformations that determine the domains, codomains, and equations. For example, a monoidal structure on $\underline{A}$ already induces functors of the form $x \mapsto x^{\otimes n}$, which are the necessary building blocks to formulate internal bialgebraic structures. In Ulmer's framework, these compound functors must be selected and carried explicitly as part of the given families of support functors, without formalizing how they systematically arise from the underlying background structure. Our syntactic framework builds from a signature pair $(\Sigma, \sigma)$ and the associated term graph $G_\sigma$ these compound functors and natural transformations. For practical applications, this syntactic shortcut is useful because the logical grammar automatically allows the constructions for the support functors that Ulmer denotes $F_{d\mu},F_{c\mu}, F_{dr},F_{cr}$ and natural transformations $\alpha_r,\beta_r$ for $\mu\in M$ and $r\in R$.

We first establish the syntactic language by introducing signature pairs $(\Sigma,\sigma)$ and bialgebraic $\sigma$-theories $T$. We then define the semantic notions in PIE-complete $2$-category $\K$: a $\Sigma$-model $M$ and construct the $0$-cells $M^\sigma$ of $\sigma$-bialgebras and $M^T$ of $T$-bialgebras. Finally, we prove the Induced Functor of Algebras Theorem (Theorem~\ref{thm:induced_morphism_between_algebras}), which establishes how a lax, oplax, or strong morphism $F\colon M\to N$ between $\Sigma$-models induces a morphism $F^T\colon M^T\to N^T$ between the corresponding objects of bialgebras in any PIE-complete $2$-category $\K$. This result unifies several disjoint theorems in categorical algebra, such as the classical lifting of lax monoidal functors to categories of internal monoids.

We leave the development of $2$-algebraic equational theories for later work. The definition of a monoid does not require that the monoidal structure is coherent. For our purposes of unifying the $1$-dimensional algebras, the structure of a $2$-algebra is more fundamental than the coherence laws the $2$-algebra might internally satisfy.

\subsection{Syntax: Signature Pairs and Theories}
To motivate the syntactic constructions in this section, we can consider the syntactic structure allowing the formulation of internal monoids in a lax monoidal setting:
$$
\Sigma = (c, \quad cc\xrightarrow{\square} c,\ ()\xrightarrow{I} c, \quad (x\square y)\square z\xrightarrow{\alpha^{xyz}}x\square (y\square z),\ x\square I\xrightarrow{\rho^x} x,\ I\square x\xrightarrow{\lambda^x} x),
$$
a relative $1$-signature describing the monoid's multiplication and unit operations
\[
\sigma = (s\colon c, \quad \mu\colon s\square s\to s, \quad u\colon I\to s),
\]
and the associative and unit equations $T$. This information is packaged within a signature pair $(\Sigma,\sigma)$ and a $\sigma$-theory $T$.

Given a function $f \colon X \to Y$ between sets, we say that the set $X$ is \emph{typed} in $Y$ and write $x \colon f(x)$ for each element $x$ in $X$.

\begin{definition}[The $2$-Signature $\Sigma$]
   A \textbf{$2$-signature} $\Sigma = (C,\mathcal{F}, \mathcal{T})$ consists of the following data:
   \begin{enumerate}
       \item A small set $C$ of \textbf{category symbols}, often denoted with $c,d,c_1,c_2,\dots$.
       \item A typed small set $\mathcal{F} \to C^* \times C$ of \textbf{functor symbols} $B \colon c_1 \cdots c_n \to d$, where $C^*$ is the free monoid over $C$.
       \item A typed set $V \to C$, with countably infinite fibers, of \textbf{variable symbols} $x,y,z,\dots$. From these variables, we recursively construct a typed set $\Term^\Sigma_V \to C$ of \textbf{$\Sigma$-terms}.\footnote{A term $t\in \Term_V^\Sigma$ is either a variable $x\colon c$ in $V$, a constant $B\colon ()\to c$ (functor symbol with empty domain) or $t = B(t_0,\ldots, t_n)\colon c$ for some previously constructed terms $t_0\colon c_0,\ldots, t_n\colon c_n$ and a functor symbol $B\colon c_0\cdots c_n\to c\in \mathcal{F}$, where $n\in \mathbb{N}$ ($0\in \mathbb{N}$).} The set of \textbf{contexts} $\mathrm{Cont}_V \subset V^*$ consists of sequences where no variable occurs more than once. A term $t \in \mathrm{Term}_V^\Sigma$ is in a context $\bar{x} \in \mathrm{Cont}_V$ if all variables appearing in $t$ occur in the sequence $\bar{x}$.
       \item A small set $\mathcal{T} \to \mathrm{Cont}_V \times \mathrm{Term}^\Sigma_V \times \mathrm{Term}^\Sigma_V$ of \textbf{transformation symbols} $\alpha^{\bar{x}} \colon t_1 \to t_2$, where $t_1, t_2$ are $\Sigma$-terms of the same type in context $\bar{x}$.
   \end{enumerate}
\end{definition}

We define the internal algebraic structures on top of this background categorical universe by introducing a relative $1$-signature.

\begin{definition}[The 1-Signature $\sigma$]
   A pair $(\Sigma,\sigma)$ is called a \textbf{signature pair} if $\Sigma = (C,\mathcal{F},\mathcal{T})$ is a $2$-signature and the \textbf{1-signature} $\sigma = (S, F)$ relative to $\Sigma$ consists of:
   \begin{enumerate}
       \item A typed small set of \textbf{sorts} $S \to C$, assigning each sort $s \in S$ a type $c_s \in C$.
       \item A typed small set of \textbf{function symbols} $F \to \mathrm{Term}_{S}^\Sigma \times \mathrm{Term}_{S}^\Sigma$. Each $f \in F$ is associated with a tuple $(t_1,t_2)$, denoted $f\colon t_1 \to t_2 \colon c_f$, where $t_1, t_2\in \Term_S^\Sigma$ are terms of type $c_f$.
   \end{enumerate}
\end{definition}

We recursively generate a graph of terms to develop equations between these operations.

\begin{definition}[The Term Graph]
   Let $(\Sigma,\sigma)$ be a signature pair. We construct the \textbf{term graph} $G_\sigma$ in the following recursive way:
   \begin{itemize}
       \item The set of vertices of $G_\sigma$ is the set $\mathrm{Term}^\Sigma_S$ of $\Sigma$-terms over the set of sorts $S$ of $\sigma$.
       \item An arrow of $G_\sigma$ is recursively defined by:
       \begin{itemize}
           \item Each function symbol $f \colon t_1 \to t_2$ of $\sigma$ forms an arrow $t_1 \to t_2$ in $G_\sigma$. 
           \item For each transformation symbol $\alpha^{\bar{x}} \colon t_1 \to t_2$ in $\Sigma$, where $\bar{x} = x_1\cdots x_n$, we set $\alpha_{t_1',\dots, t_n'} \colon t_1^{\bar{x}}(t_1',\dots, t_n') \to t_2^{\bar{x}}(t_1',\dots, t_n')$ as an arrow in $G_\sigma$ for each sequence of terms $t_i' \in \mathrm{Term}_S^\Sigma$ of type $c_{x_i}$.\footnote{The term $t_i^{\bar{x}}(t_1',\ldots, t_n')$ refers to recursively defined substitution, where $t_j'$ is substituted on $x_j$ within $t_i$ for $i = 1,2$ and $j\leq n$.}
           \item For any functor symbol $B \colon c_0\dots c_n \to d$ in $\Sigma$, a sequence of terms $t_j \colon c_j$, and an arrow $g \colon t_i \to t_i'$ in $G_\sigma$, the symbol sequence 
           $$
           B(t_0,\dots, t_{i-1}, g, t_{i+1}, \dots, t_n) \colon B(t_0,\dots, t_i, \dots, t_n) \to B(t_0,\dots, t_i', \dots, t_n)
           $$
           defines an arrow in $G_\sigma$.\footnote{We use the indexing starting from $0$ in terms $B(t_0,\ldots, t_n)$ to indicate that $B$ is not a constant symbol. If the indexing starts from $1$, it implies the possibility that $B(t_1,\ldots, t_n)$ is a constant symbol with the choice $n = 0$.}
       \end{itemize}
   \end{itemize}
\end{definition}

Using this term graph, we define equations and classify theories based on their arity and coarity behaviors.

\begin{definition}[Arity, Coarity, and Bialgebraic Theories]
   Let $(\Sigma,\sigma)$ be a signature pair. A parallel path $(p, q \colon t_1 \to t_2\colon c)$ in the term graph $G_\sigma$ is called a \textbf{$\sigma$-equation}, and a set $T$ of  $\sigma$-equations is called a \textbf{bialgebraic $\sigma$-theory}. We define the arity and coarity terms associated with $T$ as follows:
   \begin{enumerate}
       \item \textbf{Function Arity and Coarity:} For each function symbol $f \colon t_1 \to t_2$ in the $1$-signature $\sigma$, the domain term $t_1$ is a \textbf{function arity term} and the codomain term $t_2$ is a \textbf{function coarity term}.
       \item \textbf{Equational Arity and Coarity:} For each equation $(p, q \colon t_1' \to t_2')$ in the theory $T$, the domain term $t_1'$ is an \textbf{equational arity term} and the codomain term $t_2'$ is an \textbf{equational coarity term}.
       \item \textbf{Arity Terms:} The collection of all function arity terms and equational arity terms is the set of \textbf{arity terms}.
       \item \textbf{Coarity Terms:} The collection of all function coarity terms and equational coarity terms is the set of \textbf{coarity terms}.
   \end{enumerate}
   A bialgebraic $\sigma$-theory $T$ is called \textbf{algebraic} if all coarity terms consist solely of sorts $s \in S$. Similarly, the theory $T$ is called \textbf{coalgebraic} if all arity terms are sorts. 
\end{definition}

\subsection{Semantics: Models and the Object of Bialgebras}

We now define the semantics of these signatures by interpreting them as models, operations, and equations within an arbitrary PIE-complete $2$-category.

\begin{definition}[$\Sigma$-Models]
   Let $\Sigma = (C,\mathcal{F}, \mathcal{T})$ be a $2$-signature. A \textbf{$\Sigma$-model} $M$ in a $2$-category $\K$ with finite products assigns:
   \begin{enumerate}
       \item An object $M_c$ in $\K$ to each category symbol $c \in C$. For any typed sequence $\bar{x} = (x_1 \colon c_1, \dots, x_n \colon c_n)$; $c_1,\ldots, c_n\in C$; we define the product object $M_{\bar{x}} = \prod_{i=1}^n M_{c_i}$ in $\K$.
       \item A $1$-cell $M(B) \colon M_{\bar{c}} \to M_d$ in $\K$ for each functor symbol $B \colon \bar{c} \to d \in \mathcal{F}$. For any $\Sigma$-term $t \colon c$ in a context $\bar{x}$, we recursively construct the term $1$-cell $M_{\bar{x}}(t) \colon M_{\bar{x}} \to M_c$.\footnote{Consider $t\in \Term^\Sigma_V$ in context $\bar{x}$. If $t = x_i$, then $M_{\bar{x}}(t) = \pi_{x_i}\colon M_{\bar{x}}\to M_{x_i}$ and if $t = B(t_1,\ldots, t_n)$, then $M_{\bar{x}}(t) = M(B)(M_{\bar{x}}(t_1),\ldots, M_{\bar{x}}(t_n))$.}
       \item A $2$-cell $M(\alpha) \colon M_{\bar{x}}(t_1) \Rightarrow M_{\bar{x}}(t_2)$ in $\K$ for each transformation symbol $\alpha^{\bar{x}} \colon t_1 \to t_2 \in \mathcal{T}$.
   \end{enumerate}
\end{definition}

Between different $\Sigma$-models, we can define three natural notions of morphism.

\begin{definition}[Lax, Oplax, and Strong Morphisms]\label{def:laxStrongMorphisms}
   Let $\Sigma$ be a $2$-signature with $\Sigma$-models $M$ and $N$ in a $2$-category $\K$ with products. A \textbf{lax morphism} $(F,\eta) \colon M \to N$ consists of the following data:
   \begin{enumerate}
       \item A $1$-cell $F_c \colon M_c \to N_c$ for each category symbol $c$ in $\Sigma$. We define $F_{\bar{x}} = \prod_{i\leq n} F_{c_i}\colon M_{\bar{x}}\to N_{\bar{x}}$ for every typed sequence $\bar{x} = (x_1\colon c_1,\ldots, x_n\colon c_n)$, where $c_1,\ldots, c_n$ are category symbols of $\Sigma$. 
       \item A $2$-cell $\eta^B \colon N(B) \circ F_{\bar{c}} \Rightarrow F_d \circ M(B)$ for every functor symbol $B \colon \bar{c} \to d$ in $\Sigma$. We recursively extend this to define a $2$-cell $\eta_{t, \bar{x}} \colon N_{\bar{x}}(t) \circ F_{\bar{x}} \Rightarrow F_c \circ M_{\bar{x}} (t)$ for every $\Sigma$-term $t \colon c$ in a context $\bar{x} = x_1\cdots x_n$ as follows:
   $$
\eta_{t,\bar{x}} = \begin{cases}
   \id\colon \pi_i\circ F_{\bar{x}}\Rightarrow F_{c_i}\circ \pi_i, & \text{if $t = x_i$ for $i\leq n$} \\
   N(B)(N_{\bar{x}}(t_1)F_{\bar{x}},\ldots, N_{\bar{x}}(t_n)F_{\bar{x}}) \\
\quad \xRightarrow{N(B)*(\eta_{t_1,\bar{x}},\dots, \eta_{t_n, \bar{x}})} N(B)F_{\bar{c}}(M_{\bar{x}}(t_1),\dots, M_{\bar{x}}(t_n))\\
   \quad \xRightarrow{\eta^B(M_{\bar{x}}(t_1),\dots, M_{\bar{x}}(t_n))} F_c M(B)(M_{\bar{x}}(t_1),\dots, M_{\bar{x}}(t_n)) & \text{if $t = B(t_1,\dots, t_n)$}
\end{cases}
       $$
       \item For every transformation symbol $\theta^{\bar{x}} \colon t_1 \to t_2$ of type $c$ in $\Sigma$, we require that the diagram
       $$
       \begin{tikzcd}[ampersand replacement=\&]
       N_{\bar{x}}(t_1) \circ F_{\bar{x}} \arrow[r, "\eta_{t_1, \bar{x}}", Rightarrow] \arrow[d, "N(\theta)*F_{\bar{x}}"', Rightarrow] \& F_c \circ M_{\bar{x}}(t_1) \arrow[d, "F_c*M(\theta)", Rightarrow] \\
       N_{\bar{x}}(t_2) \circ F_{\bar{x}} \arrow[r, "\eta_{t_2, \bar{x}}"' , Rightarrow]                                               \& F_c \circ M_{\bar{x}}(t_2)
       \end{tikzcd}
       $$
       commutes.
   \end{enumerate}
   Dually, we define an \textbf{oplax morphism} by reversing the direction of the $2$-cells $\eta^B$. We say that $(F,\eta)$ is a \textbf{strong morphism} if it is a lax morphism where each $2$-cell $\eta_B$ is an isomorphism.
\end{definition}

\begin{proposition}
   Let $(F,\eta) \colon M \to N$ and $(G,\theta) \colon N \to P$ be lax morphisms of $\Sigma$-models. Then the composite $G \circ F$ is a lax morphism. Furthermore, oplax and strong morphisms are closed under composition.
\end{proposition}
\begin{proof}
   We sketch the case for lax morphisms. We define the components $(G \circ F)_c = G_c \circ F_c$ and define the structural $2$-cell $\phi^B$ for each functor symbol $B \colon \bar{c} \to d$ as the composite:
   $$
   \phi^B \colon P(B) \circ (GF)_{\bar{c}} \xRightarrow{\theta^B*F_{\bar{c}}} G_d \circ N(B) \circ F_{\bar{c}} \xRightarrow{G_d * \eta^B} (GF)_d \circ M(B)
   $$
   By induction, for any $\Sigma$-term $t \colon c$ in a context $\bar{x}$, we obtain the factorization:
   $$
   \phi_{t, \bar{x}} \colon P_{\bar{x}}(t) \circ (GF)_{\bar{x}} \xRightarrow{\theta_{t, \bar{x}}*F_{\bar{x}}} G_c \circ N_{\bar{x}}(t) \circ F_{\bar{x}} \xRightarrow{G_c*\eta_{t, \bar{x}}} (GF)_c \circ M_{\bar{x}}(t)
   $$
   Consider a transformation symbol $\alpha^{\bar{x}} \colon t_1 \to t_2\colon c$ in $\Sigma$. Now the diagram
   $$
   \begin{tikzcd}[column sep = 2cm, ampersand replacement=\&]
   P_{\bar{x}}(t_1) \circ (GF)_{\bar{x}} \arrow[d, "P(\alpha)*(GF)_{\bar{x}}" description, Rightarrow] \arrow[rr, "\phi_{t_1, \bar{x}}", Rightarrow, bend left] \arrow[r, "\theta_{t_1, \bar{x}}*F_{\bar{x}}" description, Rightarrow] \& G_c \circ N_{\bar{x}}(t_1) \circ F_{\bar{x}} \arrow[r, "G_c*\eta_{t_1, \bar{x}}" description, Rightarrow] \arrow[d, "G_c*N(\alpha)*F_{\bar{x}}" description, Rightarrow] \& (GF)_c \circ M_{\bar{x}}(t_1) \arrow[d, "(GF)_c*M(\alpha)" description, Rightarrow] \\
   P_{\bar{x}}(t_2) \circ (GF)_{\bar{x}} \arrow[r, "\theta_{t_2, \bar{x}}*F_{\bar{x}}"' description, Rightarrow] \arrow[rr, "\phi_{t_2, \bar{x}}"', Rightarrow, bend right]                                        \& G_c \circ N_{\bar{x}}(t_2) \circ F_{\bar{x}} \arrow[r, "G_c*\eta_{t_2, \bar{x}}"' description, Rightarrow]                                                                \& (GF)_c \circ M_{\bar{x}}(t_2)
   \end{tikzcd}
   $$
   commutes. Thus, the composite $G \circ F$ is a lax morphism of $\Sigma$-models.
\end{proof}

We obtain the $2$-category of $\Sigma$-models with lax morphisms, and its sub-$2$-category of strong morphisms, by defining the $2$-cells. A transformation $\gamma \colon (F,\eta_F) \Rightarrow (G, \eta_G) \colon M \to N$ of lax $\Sigma$-morphisms consists of $2$-cells $\gamma_c \colon F_c \Rightarrow G_c$ in $\K$ such that the diagram:
$$
\begin{tikzcd}[ampersand replacement=\&]
N(B) \circ F_{\bar{c}} \arrow[d, "N(B)*\gamma_{\bar{c}}" description, Rightarrow] \arrow[r, "\eta^B_F", Rightarrow] \& F_d \circ M(B) \arrow[d, "\gamma_d*M(B)", Rightarrow] \\
N(B) \circ G_{\bar{c}} \arrow[r, "\eta_G^B"', Rightarrow]                                                           \& G_d \circ M(B)
\end{tikzcd}
$$
commutes for each functor symbol $B\colon \bar{c}\to d$ in $\Sigma$. Similarly, one attains the $2$-category of $\Sigma$-models with oplax morphisms and previously defined $2$-cells.

\begin{definition}[The Object of Operations $M^\sigma$]
   Consider a signature pair $(\Sigma,\sigma)$ with $S$ as the set of $\sigma$-sorts. Consider a $\Sigma$-model $M$ in a $2$-category $\K$ with small products and inserters. With products, we define $M_0 = \prod_{s \in S} M_{c_s}$. We define the object $M^\sigma$ of operations as the inserter
   $$
   \begin{tikzcd}
M^\sigma \arrow[d, "u_M"'] \arrow[r, "u_M"]                                    & M_0 \arrow[d, "(M(t_2))_f"]            \\
M_0 \arrow[r, "(M(t_1))_f"'] \arrow[ru, "\phi_M" description, Rightarrow] & \prod_{f} M_c
\end{tikzcd}
   $$
   where the indexing is over $\sigma$-function symbols $f\colon t_1\to t_2\colon c$.\footnote{For a term $t\in\Term_S^\Sigma$, we define $M(t)\colon M_0\to M_c$ recursively: $M(t) = \pi_s\colon M_0\to M_{c_s}$ if $t = s$ is a sort and $M(t) = M(B)(M(t_1),\ldots, M(t_n))$, if $t = B(t_1,\ldots, t_n)$.} For a term $t\colon c\in \Term_S^\Sigma$, we denote $M^\sigma(t) = M(t)\circ u_M\colon M^\sigma\to M_c$. If $\K = \CAT$, we call an object of $M^\sigma$ a \textbf{$\sigma$-model}.
\end{definition}
Note that a $\sigma$-model $m$ in $M$ consists of a choice of an object $m_s$ and a morphism $m(f)\colon M(t_1)((m_s)_s)\to M(t_2)((m_s)_s)$ for each $\sigma$-sort $s$ and a $\sigma$-function symbol $f\colon t_1\to t_2$. 

\begin{definition}[Path Evaluation and the Object of Bialgebras $M^T$]
   Let $(\Sigma,\sigma)$ be a signature pair with a $\Sigma$-model $M$ in a PIE-complete $2$-category $\K$. We define a $2$-cell $\overline{p} \colon M^\sigma(t_1) \Rightarrow M^\sigma(t_2) \colon M^\sigma \to M_c$ recursively for each arrow $p \colon t_1 \to t_2 \colon c$ in the term graph $G_\sigma$:
   {\footnotesize
   $$
   \overline{p}_M = \begin{cases}
       \pi_f * \phi \colon M^\sigma(t_1) \Rightarrow M^\sigma(t_2), & \text{if $p = f \colon t_1 \to t_2$ is a function symbol in $\sigma$,} \\
       M(\alpha) * (M^\sigma(t_1'), \dots, M^{\sigma}(t_n')), & \text{if $p = \alpha_{t_1', \dots, t_n'}, \alpha^{\bar{x}}\colon t_1\to t_2$ a $\Sigma$-transformation symbol,} \\
       M(B) * (M^\sigma(t_1'), \dots, \bar{q}_M, \dots, M^\sigma(t_n')), & \text{if $p = B(t_1', \dots, q, \dots, t_n')$, $B$ is a $\Sigma$-functor symbol.}
   \end{cases}
   $$
   }
   This extends via vertical composition to all paths $p$ in $G_\sigma$.
   
   For a bialgebraic $\sigma$-theory $T$, the \textbf{object of $T$-bialgebras} $M^T$ is defined via the equifier
   $$
       \begin{tikzcd}[ampersand replacement=\&]
{M^T} \& {M^\sigma} \&\& {\prod_{(p,q)} M_c}
\arrow["{e_M}", from=1-1, to=1-2]
\arrow[""{name=0, anchor=center, inner sep=0}, "{(M^\sigma(t_1))}", shift left=3, from=1-2, to=1-4]
\arrow[""{name=1, anchor=center, inner sep=0}, "{(M^\sigma(t_2))}"', shift right=3, from=1-2, to=1-4]
\arrow["{(\bar{p}_M)}"', shift right=3, Rightarrow, from=0, to=1]
\arrow["{(\bar{q}_M)}", shift left=3, Rightarrow, from=0, to=1]
\end{tikzcd}
   $$
   where the indexing is over $(p,q\colon t_1\to t_2\colon c)\in T$. If $\K = \CAT$, we call an object of $M^T$ a \textbf{$T$-bialgebra}.
\end{definition}
The first part of the following lemma justifies that the transformation $\overline{\alpha_{t_1',\ldots, t_n'}}_M$ in the previous definition has the correct domain and codomain. For a signature pair
$(\Sigma,\sigma = (S,F))$ and a lax morphism $(F,\eta)\colon M\to N$ of $\Sigma$-models in $\K$ with products, we define 
$$
F_0\coloneqq  \prod_{s\in S} F_{c_s}\colon M_0\to N_0,\quad \eta_t\colon N(t)F_0\Rightarrow F_cM(t)\colon M_0\to N_c
$$
for $t\colon c\in \Term_S^\Sigma$ almost the same as in Definition \ref{def:laxStrongMorphisms}(2); one just needs to substitute $0$ on $\bar{x}$.\footnote{Note that we have $1$-cells $M_{\bar{x}}(t)\colon M_{\bar{x}}\to M_c$ and $M(t')\colon M_0\to M_c$ for $t\colon c\in \Term_V^\Sigma$ in context $\bar{x}$ and $t'\in \Term_S^\Sigma$.}

\begin{lemma}\label{lem:substitution_evaluated}
   Let $(F,\eta) \colon M \to N$ be a lax morphism of $\Sigma$-models in $\K$. Let us denote by $V$ and $S$ the sets of variable symbols of $\Sigma$ and sorts of $S$, respectively. Let $t\in \Term_V^\Sigma$ be a term of type $c$ in a context $\bar{x} = x_1\cdots x_n$, where $x_i\colon c_i$ for category symbols $c,c_i,i\leq n$. Let $t_1\colon c_1,\ldots, t_n\colon c_n\in \Term_S^\Sigma$. Then we have the following factorizations
   \begin{align}
           M(t^{\bar{x}}(t_1,\ldots, t_n)) &\colon M_0\xrightarrow{(M(t_1),\ldots, M(t_n))}M_{\bar{x}}\xrightarrow{M_{\bar{x}}(t)} M_c\label{eq:0}\\
           \eta_{t^{\bar{x}}(t_1,\ldots, t_n)} &\colon N(t^{\bar{x}}(t_1,\ldots,t_n))F_0\xRightarrow{ N_{\bar{x}}(t)(\eta_{t_1},\ldots, \eta_{t_n})}N_{\bar{x}}(t)F_{\bar{x}}(M(t_1),\ldots, M(t_n))
           \nonumber\\
           &\quad \xRightarrow{\eta_{t,\bar{x}}(M(t_1),\ldots, M(t_n))}F_cM(t^{\bar{x}}(t_1,\ldots, t_n))\label{eq:1}
   \end{align}\qedhere
\end{lemma}
\begin{proof}
We prove the claim by induction on the structure of the term $t$.

\noindent\textbf{Case 1:} Assume $t = x_i$. Now $t^{\bar{x}}(t_1,\ldots, t_n) = t_i$. By definition, we have that $x_i^{\bar{x}}(t_1, \dots, t_n) = t_i$ and now $M(t_i) = \pi_i(M(t_1),\ldots, M(t_n))$ and $\eta_{t_i} = id\bullet \pi_i(\eta_{t_1},\ldots, \eta_{t_n})$. Thus the equations hold.


\noindent\textbf{Case 2:} Assume $t = B(t_1',\ldots, t_k')$, where the equations hold for $t_i'$ for $i\leq k$ and $B$ is a functor symbol in $\Sigma$. 
The substitution on $t$ then yields:
\[
t^{\bar{x}}(t_1, \dots, t_n) = B(t_1'^{\bar{x}}(t_1,\ldots, t_n),\ldots, t_k'^{\bar{x}}(t_1,\ldots, t_n))= B(s_1, \dots, s_k),
\]
where $s_i = t_i'^{\bar{x}}(t_1,\ldots, t_n)$ for $i\leq k$. By the induction hypothesis $M(s_i) = M_{\bar{x}}(t_i')\bar{M}$, where $\bar{M}\coloneqq (M(t_1),\ldots, M(t_n))$ for $i\leq k$.
To verify Equation \eqref{eq:0}, consider:
\begin{align*}
   &M(t^{\bar{x}}(t_1, \dots, t_n)) \\
   &= M(B(s_1, \dots, s_k)) \\
   &= M(B)(M(s_1), \dots, M(s_k)) \tag{Definition of $M(B(s_1,\dots, s_k))$} \\
   &= M(B) (M_{\bar{x}}(t'_1)\bar{M}, \dots, M_{\bar{x}}(t'_k)\bar{M}) \tag{Induction hypothesis \eqref{eq:0} on $t'_i$} \\
   &= M(B) (M_{\bar{x}}(t'_1), \dots, M_{\bar{x}}(t'_k)) \bar{M} \\
   &= M_{\bar{x}}(B(t'_1, \dots, t'_k))\bar{M} \\
   &= M_{\bar{x}}(t)(M(t_1), \dots, M(t_n))
\end{align*}
This proves equation \eqref{eq:0} for the inductive step. Denote $\bar{\eta}\coloneqq (\eta_{t_1},\ldots, \eta_{t_n})$. For the equation \eqref{eq:1}, consider:
{
\begin{align*}
   &\eta_{t^{\bar{x}}(t_1,\ldots, t_n)}\\
   &= \eta_{B(s_1,\ldots, s_k)}\\
   &= \eta^B (M(s_1),\ldots, M(s_k)) \bullet N(B)(\eta_{s_1},\ldots, \eta_{s_k}) \tag{Definition of $\eta_{B(s_1,\ldots, s_k)}$}\\
   &= \eta^B(M_{\bar{x}}(t_1')\bar{M},\ldots, M_{\bar{x}}(t_k')\bar{M})\bullet\\
   & N(B)(\eta_{t_1',\bar{x}}\bar{M} \bullet N_{\bar{x}}(t_1')\bar{\eta},\ldots, \eta_{t_k',\bar{x}}\bar{M}\bullet N_{\bar{x}}(t_k')\bar{\eta}) \tag{Induction hypothesis on $t_i'$}\\
   &=\eta^B(M_{\bar{x}}(t_1')\bar{M},\ldots, M_{\bar{x}}(t_k')\bar{M})\bullet
   N(B)(\eta_{t_1',\bar{x}}\bar{M},\ldots, \eta_{t_k',\bar{x}}\bar{M})\bullet\\
   & N(B)(N_{\bar{x}}(t_1')\bar{\eta},\ldots, N_{\bar{x}}(t_k')\bar{\eta}) \tag{Composition componentwise}\\
   &= (\eta^B(M_{\bar{x}}(t_1'),\ldots, M_{\bar{x}}(t_k'))\bullet
   N(B)(\eta_{t_1',\bar{x}},\ldots, \eta_{t_k',\bar{x}}))\bar{M} \bullet
    N(B)(N_{\bar{x}}(t_1'),\ldots, N_{\bar{x}}(t_k'))\bar{\eta} \\
   &= \eta_{B(t_1',\ldots, t_k'),\bar{x}}\bar{M}\bullet N_{\bar{x}}(B(t_1',\ldots, t_k'))\bar{\eta}\tag{Definitions of $\eta_{B(t_1',\ldots, t_k')}$ and $N_{\bar{x}}(B(t_1',\ldots, t_k'))$}\\
   &=\eta_{t,\bar{x}}(M(t_1),\ldots, M(t_n))\bullet N_{\bar{x}}(t)(\eta_{t_1},\ldots, \eta_{t_n})\tag{$t = B(t_1',\ldots, t_k')$}
\end{align*}
}
\end{proof}

\subsection{Examples}

\begin{examples}[Examples of $2$-Signatures]\label{ex:$2$-signatures}
   We define the following standard $2$-signatures representing background categorical structures. 
   {\small
   \begin{align*}
       \Sigma_{\text{Cart}}    
    &= (c,\quad \square\colon cc\to c,\ I\colon ()\to c,\ S\colon c\to c,
    \quad x\xrightarrow{\delta^x}x\square x, \ x\xleftarrow{p_1^{xy}}x\square y\xrightarrow{p_2^{xy}}y,\ x\xrightarrow{!^x}I)\\
       \Sigma_{\text{Mon}}     
       &= (c,\quad \square\colon cc\to c,\  I\colon ()\to c,\quad (x\square y)\square z\rightleftarrows x\square (y\square z),\ x\square I\rightleftarrows x\leftrightarrows I\square x)\\
       \Sigma_{\text{SMon}}
       &= (c,\quad \square\colon cc\to c,\ I\colon ()\to c,\quad (x\square y)\square z\rightleftarrows x\square (y\square z),\  x\square I\rightleftarrows x\leftrightarrows I\square x,\ x\square y\to y\square x)\\
       \Sigma_{\text{Adj}}
       &= (c,d,\quad L\colon c\to d,\ R\colon d\to c,\quad \eta^x\colon x\to R(L(x)),\ \varepsilon_y\colon L(R(y))\to y)\\
       \Sigma_{\text{Mnd}}
       &= (c,\quad T\colon c\to c,\quad \mu^x\colon T(T(x))\to T(x),\ \eta^x\colon x\to T(x))\\
       \Sigma_{\text{Distr}}
       &= (c,\quad T,S\colon c\to c,\quad  \mu_T^x,\ \eta_T^x,\ \mu_S^x,\ \eta_S^x,\ \gamma^x\colon S(T(x))\to T(S(x))) 
   \end{align*}
   }
\end{examples}

\begin{examples}\label{exs:theories}
   We provide examples for the theories of monoids, Hopf algebras and fields:
   \begin{enumerate}
       \item \textbf{Monoids:} Consider the signature pair $(\Sigma, \sigma)$ defined as follows:
       \begin{align*}
       \Sigma   
       &= (c,\quad \square\colon cc\to c, I\colon ()\to c,\quad (x\square y)\square z\xrightarrow{\alpha^{xyz}} x\square (y\square z), x\square I\xrightarrow{\rho^x} x\xleftarrow{\lambda^x} I\square x)\\
       \sigma
       &= (s\colon c,\quad \mu\colon s\square s\to s, u\colon I\to s)
   \end{align*}
   For convenience, we write $x\square y$ for $\square(x,y)$. Consider the following parallel paths in the term graph $G_{\sigma}$ defining the theory of monoids $T$:
   $$
   \begin{tikzcd}[column sep=0.8 cm]
   (s \square s) \square s \arrow[d, "\mu \square s"'] \arrow[r, "\alpha_{s,s,s}"] & s \square (s \square s) \arrow[r, "s \square \mu"] & s \square s \arrow[d, "\mu"] & s \square I \arrow[r, "s \square u"] \arrow[rd, "\rho_s"'] & s \square s \arrow[d, "\mu"] & s \square s \arrow[d, "\mu"'] & I \square s \arrow[l, "u \square s"'] \arrow[ld, "\lambda_s"] \\
   s \square s \arrow[rr, "\mu"'] & & s & & s & s & 
   \end{tikzcd}
   $$
   Let $M$ be a $\Sigma$-model in $\CAT$, a monoidal category, for instance. We denote $\otimes\coloneqq M(\square)$ and denote $I$ for $M(I)(*)$, where $*$ is the unique element of the terminal category. Note that 
   $$
   M^\sigma = \Ins((M(s\square s),M(I)), (M(s),M(s))),
   $$
   where $M(s\square s)$, $M(s)$ and $M(I)$ are the functors 
   $$
   x\mapsto x\otimes x, \quad x, \quad I\colon M_c\to M_c\text{, respectively.} 
   $$
   Thus $M^\sigma$ consists of objects $(x,m\colon x\otimes x\to x, e\colon I\to x)$, and $M^\sigma$ can be understood as the category of pointed magmas of $M$. Consider the forgetful functor $U_M\colon M^\sigma\to M_c$. To extract the subcategory of $M^\sigma$ of internal monoids, we need to consider a suitable full subcategory $M^T$.

   By definition, the category $M^T$ is the equifier of the parallel pair of natural transformations:
   \begin{align*}
   &\left(\overline{\mu \circ (\mu \square s)}_M, \,\, \overline{\rho_s}_M, \,\, \overline{\lambda_s}_M\right), \,\, \left(\overline{\mu \circ (s \square \mu) \circ \alpha_{s,s,s}}_M, \,\, \overline{\mu \circ (s \square u)}_M, \,\, \overline{\mu \circ (u \square s)}_M\right)\colon \\
   &\left(M((s \square s)\square s)U_M, \,\, M(s \square I)U_M, \,\, M(I \square s)U_M\right) \Rightarrow (U_M, \,\, U_M, \,\, U_M)\colon M^\sigma\to M_c^3
   \end{align*}
   This means an object $(x, m, e)$ of $M^\sigma$ is in $M^T$ if and only if 
   {\small
   $$
   \left(\overline{\mu \circ (\mu \square s)}_M, \,\, \overline{\rho_s}_M, \,\, \overline{\lambda_s}_M\right)_{(x,m,e)} = \left(\overline{\mu \circ (s \square \mu) \circ \alpha_{s,s,s}}_M, \,\, \overline{\mu \circ (s \square u)}_M, \,\, \overline{\mu \circ (u \square s)}_M\right)_{(x,m,e)},
   $$
   which is exactly to say that the following three equations hold in $M_c$:
   \begin{align*}
       m \circ (m \otimes x) &= m \circ (x \otimes m) \circ M(\alpha)_{x,x,x} \\
       m \circ (x \otimes e) &= M(\rho)_x \\
       m \circ (e \otimes x) &= M(\lambda)_x
   \end{align*}
   }
   These are precisely the classical associativity and unit laws for a monoid in $M$.

\item \textbf{Hopf Algebras:}
    To define internal Hopf algebras, we work relative to a skew-monoidal $2$-signature with a chosen interchange law:
    {\small
    \begin{align*}
        \Sigma&=(c, \quad \square \colon cc \to c,\ I \colon () \to c, \\
        &\quad (x \square y) \square z\xrightarrow{\alpha^{xyz}}  x \square (y \square z),\ I \square x \xrightarrow{\lambda^x} x \xrightarrow{\rho^x} x \square I,\ (x \square y) \square (z \square v) \xrightarrow{\chi^{xyzv}} (x \square z) \square (y \square v))
    \end{align*}
    }
    where $\chi^{xyzv}$ represents the interchange law, which is often uniquely constructed via the coherence of symmetric monoidal categories \cite{MacLane1998}.

    We define the relative $1$-signature:
    $$
    \sigma_{\text{Hopf}} \coloneqq (s\colon c, \quad \mu\colon s \square s \to s, \ u\colon I \to s, \ \delta\colon s \to s \square s, \ \varepsilon\colon s \to I, \ S\colon s \to s)
    $$
    The theory $T_{\text{Hopf}}$ consists of the following equations in the term graph $G_{\sigma_{\text{Hopf}}}$, presented as commutative diagrams where each face represents a parallel pair of paths:\footnote{We use $p_1\square p_2$ for the path $(p_1\square t_2')\circ (t_1\square p_2)$, where $p_i\colon t_i\to t_i'$ is an arrow in the term graph $G_{\sigma_{\mathrm{Hopf}}}$ for $i = 1,2$.}
    {\small $$
\begin{tikzcd}[column sep = 0.7 cm]
(s\square s)\square s \arrow[d, "s\square \mu"'] \arrow[r, "{\alpha_{s,s,s}}"] & s\square (s\square s) \arrow[r, "s\square \mu"]              & s\square s \arrow[d, "\mu"]                                         & I\square s \arrow[rd, "\lambda_s"'] \arrow[r, "u\square s"]              & s\square s \arrow[d, "\mu" description]                                         & s\square I \arrow[l, "s\square u"']    &                        \\
s\square s \arrow[rr, "\mu"']                                                  &                                                              & s                                                                   &                                                                          & s                                                                               & s \arrow[l, "="'] \arrow[u, "\rho_s"'] &                        \\
s \arrow[d, "\delta"'] \arrow[r, "\delta"]                                     & s\square s \arrow[r, "\delta\square s"]                      & (s\square s)\square s \arrow[d, "{\alpha_{s,s,s}}"]                 & I\square s \arrow[d, "\lambda_s"']                                       & s\square s \arrow[l, "\varepsilon\square s"'] \arrow[r, "s\square \varepsilon"] & s\square I                             &                        \\
s\square s \arrow[rr, "s\square \delta"']                                      &                                                              & s\square(s\square s)                                                & s                                                                        & s \arrow[l, "="'] \arrow[u, "\delta" description] \arrow[ru, "\rho_s"']         &                                        &                        \\
(s\square s)\square (s\square s) \arrow[r, "{\chi_{s,s,s,s}}"]                 & (s\square s)\square (s\square s) \arrow[r, "\mu\square \mu"] & s\square s                                                          & I\square I \arrow[r, "\lambda_I"]                                        & I                                                                               & I\square I \arrow[r, "u\square u"]     & s\square s             \\
s\square s \arrow[u, "\delta\square \delta"] \arrow[rr, "\mu"']                &                                                              & s \arrow[u, "\delta"']                                              & s\square s \arrow[r, "\mu"'] \arrow[u, "\varepsilon\square \varepsilon"] & s \arrow[u, "\varepsilon"']                                                     & I \arrow[r, "u"'] \arrow[u, "\rho_I"]  & s \arrow[u, "\delta"'] \\
                                                                               & I                                                            & s\square s \arrow[rr, "s\square S"]                                 &                                                                          & s\square s \arrow[d, "\mu"]                                                     &                                        &                        \\
I \arrow[r, "u"'] \arrow[ru, "="]                                              & s \arrow[u, "\varepsilon"']                                  & s \arrow[r, "\varepsilon"] \arrow[u, "\delta"] \arrow[d, "\delta"'] & I \arrow[r, "u"]                                                         & s                                                                               &                                        &                        \\
                                                                               &                                                              & s\square s \arrow[rr, "S\square s"']                                &                                                                          & s\square s \arrow[u, "\mu"']                                                    &                                        &                       
\end{tikzcd}
$$
}

    
    Evaluating this theory in $\mathbf{CAT}$ over a symmetric monoidal category $M$ yields precisely the classical category $\mathbf{Hopf}_M$ of internal Hopf algebras.

\item \textbf{Fields:}
    The category of fields can be characterized starting from the category of commutative groups. Specifically, a field $k$ corresponds to a commutative group $k^*$ with a commutative ring structure on $k^*+1$ whose multiplication extends that of $k^*$.

    To formalize this, we consider a signature pair $(\Sigma, \sigma)$:
    \begin{align*}
    \Sigma
    &\coloneqq (c,\quad \square\colon cc\to c,\ I\colon ()\to c,\ S\colon c\to c,\\
    &\quad\quad x\xrightarrow{\delta^x}x\square x, \ x\xleftarrow{p_1^{xy}}x\square y\xrightarrow{p_2^{xy}}y,\ x\xrightarrow{!^x}I,\ x\xrightarrow{\eta^x}Sx)\\
    \sigma 
    &\coloneqq (s \colon c,\quad m\colon s\square s\to s,\ i\colon s\to s,\ u\colon I\to s,\\
    &\quad\quad m_1,m_2\colon Ss\square Ss\to Ss,\ \nu\colon Ss\to Ss,\iota\colon I\to Ss)
    \end{align*}
    To define the bialgebraic theory of fields, we introduce the theory $T$ containing precisely the following equations in the term graph $G_\sigma$:
    \begin{itemize}
        \item The equations stating that $(s, m, i, u)$ satisfies the equations of a commutative group.
        \item The equations stating that $(Ss, m_1, \nu, \iota, m_2, \eta_s \circ u)$ determines a commutative ring, where $(Ss, m_1, \nu, \iota)$ and $(Ss, m_2, \eta_s \circ u)$ satisfy the additive and the multiplicative equations, respectively.
        \item An equation stating that $\eta_s \colon s \to Ss$ determines a morphism of magmas from $(s,m)$ to $(Ss,m_2)$:
        $$
        m_2 \circ (\eta_s \square \eta_s) = \eta_s \circ m \colon s \square s \to Ss
        $$
    \end{itemize}
    We choose the categorical $\Sigma$-model $M = (\Set, \times, 1, (-)+1)$ where $M(\eta)$ is the unit of the maybe monad $(-)+1$. Now, $M^T$ is the category of fields.\footnote{One could take a bifunctor symbol $\boxplus$ with the codiagonal and injections instead of $S, \eta$ and $\iota$. This would be more advantageous in the sense that the coproduct can be defined as the left adjoint to the diagonal $1$-cell. This allows an essentially unique formulation for fields if we were to develop the notion of $2$-equational $\Sigma$-theories.}

\item \textbf{Generalizing Lax Limits:}
    Consider an arbitrary $2$-signature $\Sigma = (C,\mathcal{F},\mathcal{T})$. We construct an associated algebraic theory $T$ for $\Sigma$ which generalizes the notion of a lax limit of a strict $2$-functor to the cartesian (multivariable) setting. Specifically, this construction unifies the theories of monoids, commutative monoids, and monad algebras.
    
    We define the $1$-signature $\sigma_{\text{Lax}}$ relative to $\Sigma$:
    $$
    \sigma_{\text{Lax}} \coloneqq \left(s_c\colon c, \quad f_B\colon B(s_{c_1},\dots, s_{c_n})\to s_d\right)_{c\in C,B\colon c_1\cdots c_n\to d\in\mathcal{F}}.
    $$
    To be more precise, the notation above refers to the fact that the set $S$ of sorts is in a bijective correspondence with $C$ and $F$ with $\mathcal{F}$. We recursively define for each term $t \colon c \in \Term_S^\Sigma$ a path $f_t \colon t \to s_c$ in the term graph $G_{\sigma_{\text{Lax}}}$:
    {\small
    $$
    f_t\coloneqq \begin{cases}
        ()_{s}\colon s\to s, & \text{if $t = s\in S$,}\\
        B(t_1, \dots, t_n) \xrightarrow{B(f_{t_1}, \dots, f_{t_n})} B(s_{c_1}, \dots, s_{c_n}) \xrightarrow{f_B} s_c, & \text{if $t = B(t_1,\dots, t_n)$, $t_i\colon c_i$ for $i\leq n$,}
    \end{cases}
    $$
    }
    where we utilize the convention that for any paths $p^i = (p^i_{m_i},\dots, p^i_1)\colon t_i'\to t_i''$ in $G_\sigma$ (for $i\leq n$) and a $\Sigma$-functor symbol $B$, the term $B(p^1, \dots, p^n)$ denotes the composition:
    \begin{align*}
        B(p^1,\dots, p^n) &\coloneqq B(p^1_{m_1},t_2'',\dots, t_n'')\circ\dots\circ B(p^1_1,t_2'',\dots, t_n'') \circ \dots \\
        &\quad \dots \circ B(t_1',\dots, t_{n-1}',p^n_{m_n})\circ\dots\circ B(t_1',\dots, t_{n-1}', p^n_1).
    \end{align*}
    The theory $T_{\text{Lax}}$ consists of the following parallel paths:
    \[
    \begin{tikzcd}
    t_1^{\bar{x}}(\bar{s}) \arrow[r, "\alpha_{\bar{s}}"] \arrow[rd, "f_{t_1^{\bar{x}}(\bar{s})}"'] & t_2^{\bar{x}}(\bar{s}) \arrow[d, "f_{t_2^{\bar{x}}(\bar{s})}"] \\
                                                                                                   & s_c                                                           
    \end{tikzcd}
    \]
    for each $\alpha^{\bar{x}}\colon t_1\to t_2\in \mathcal{T}$, where $\bar{s} = s_{c_1}\cdots s_{c_n}$ and $x_i\colon c_i$ for $i\leq n$. 
    
    If $\Sigma$ is the signature for a monad, a monoidal category or a symmetric monoidal category, then the associated algebraic theory is of the monad algebras, monoids and commutative monoids, respectively.
       \end{enumerate}

\end{examples}
\subsection{Induced Functor of Algebras}

In this subsection, we fix a PIE-complete $2$-category $\K$, a signature pair $(\Sigma,\sigma)$ and a bialgebraic $\sigma$-theory $T$. Our mission is to find conditions so that a morphism $F\colon M\to N$ of $\Sigma$-models in $\K$ induces a morphism $F\colon M^T\to N^T$ between the objects of $T$-bialgebras.

\begin{lemma}\label{lem:morphism_between_inserters}
   Let $(F,\eta) \colon M \to N$ be a lax morphism of $\Sigma$-models in $\K$. Assume that $\eta_t\colon N(t)F_0\Rightarrow F_c M(t)\colon M_0\to N_c$ is an isomorphism for each function coarity term $t\colon c$. Then the following holds:
   \begin{enumerate}
       \item There exists a unique morphism $F^\sigma\colon M^\sigma\to N^\sigma$ satisfying the equations
       \begin{align}
           u_N \circ F^\sigma 
           &= F_0 \circ u_M \nonumber \\
           \phi_N * F^\sigma
           &= (\eta_{t_2}^{-1}*u_M)_f \bullet ((\prod_{f\colon t_1\to t_2\colon c} F_{c}) * \phi_M) \bullet (\eta_{t_1}*u_M)_f\label{eq:2}
   \end{align}
       \item For every path $p\colon t_1\to t_2$ in the term graph $G_\sigma$, the following diagram commutes:
       $$
\begin{tikzcd}
N(t_1)F_0u_M \arrow[r, "\eta_{t_1}*u_M", Rightarrow] \arrow[d, "\bar{p}_N*F^\sigma"', Rightarrow] & F_cM(t_1)u_M \arrow[d, "F_c*\bar{p}_M", Rightarrow] \\
N(t_2)F_0u_M \arrow[r, "\eta_{t_2}*u_M"', Rightarrow]                                             & F_cM(t_2)u_M                                       
\end{tikzcd}
$$
In other words, 
\begin{equation}\label{eq:3}
\eta_{t_2}u_M\bullet \bar{p}_NF^\sigma = F_c\bar{p}_M\bullet \eta_{t_1}u_M
\end{equation}

   \end{enumerate}
\end{lemma}
\begin{proof}
   The first universal property of the inserter $N^\sigma$ induces the unique morphism $F^\sigma$:
   $$
   \begin{tikzcd}[row sep = 1 cm]
M^\sigma \arrow[d, "u_M"'] \arrow[r, "u_M"]                                                 & M_0 \arrow[r, "F_0"] \arrow[d, "(M(t_2))_f" description]                                                                              & N_0 \arrow[dd, "(N(t_2))_f"] \\
M_0 \arrow[d, "F_0"'] \arrow[r, "(M(t_1))_f"'] \arrow[ru, "\phi_M" description, Rightarrow] & \prod_{f\colon t_1\to t_2\colon c}M_c \arrow[rd, "\prod_f F_c" description] \arrow[ru, "(\eta_{t_2}^{-1})_f" description, Rightarrow] &                              \\
N_0 \arrow[rr, "(N(t_1))_f"'] \arrow[ru, "(\eta_{t_1})_f" description, Rightarrow]          &                                                                                                                                       & \prod_f N_c                 
\end{tikzcd}
   $$
   We prove by structural induction on arrows $p$ in $G_\sigma$ that the equation $\eqref{eq:3}$ holds:
       
       \noindent\textbf{Case 1:} Assume $p = f$ is a $\sigma$-function symbol. Whiskering the equation $\eqref{eq:2}$ with $\pi_f$ from the left yields the equation \eqref{eq:3}.

   \noindent\textbf{Case 2:} Assume $p = \alpha_{t_1',\ldots, t_n'}$, where $\alpha^{\bar{x}}\colon t_1''\to t_2''$ is a transformation symbol in $\Sigma$ and $t_i = t_i''^{\bar{x}}(t_1',\ldots, t_n')$ for $i = 1,2$.
   Now
     {\small \begin{align*}
       &\eta_{t_2}u_M\bullet \bar{p}_NF^\sigma\\
       &= \eta_{t_2}u_M\bullet N(\alpha)(N^\sigma(t_1'),\ldots, N^\sigma(t_n'))F^\sigma\tag{Definition of $\bar{p}_N $}\\
       &= \eta_{t_2}u_M\bullet N(\alpha)(N(t_1')F_0u_M, \ldots, N(t_n')F_0u_M)\tag{$N^\sigma(t_i')F^\sigma = N(t_i')u_N F^\sigma = N(t_i')F_0u_M$}\\
       &= (\eta_{t_2}\bullet N(\alpha)(N(t_1')F_0,\dots, N(t_n')F_0))u_M \tag{Whiskering}\\
       &= (\eta_{t_2'',\bar{x}}(M(t_1'),\ldots, M(t_n'))\bullet N_{\bar{x}}(t_2'')(\eta_{t_1'},\ldots, \eta_{t_n'})\bullet N(\alpha)(N(t_1')F_0,\dots, N(t_n')F_0))u_M \tag{\ref{eq:1}}\\
       &= (\eta_{t_2'',\bar{x}}(M(t_1'),\ldots, M(t_n'))\bullet N(\alpha)F_{\bar{x}}(M(t_1'),\dots, M(t_n'))\bullet N_{\bar{x}}(t_1'')(\eta_{t_1'},\ldots, \eta_{t_n'}))u_M\tag{Interchange}\\
       &= ((\eta_{t_2'',\bar{x}}\bullet N(\alpha)F_{\bar{x}})(M(t_1'),\ldots, M(t_n'))\bullet N_{\bar{x}}(t_1'')(\eta_{t_1'},\ldots, \eta_{t_n'}))u_M\tag{Whiskering}\\
       &= ((F_cM(\alpha)\bullet\eta_{t_1'',\bar{x}})(M(t_1'),\ldots, M(t_n'))\bullet N_{\bar{x}}(t_1'')(\eta_{t_1'},\ldots, \eta_{t_n'}))u_M\tag{Definition of lax morphism}\\
       &= (F_cM(\alpha)(M(t_1'),\ldots, M(t_n'))\bullet \eta_{t_1'',\bar{x}}(M(t_1'),\ldots, M(t_n'))\bullet N_{\bar{x}}(t_1'')(\eta_{t_1'},\ldots, \eta_{t_n'}))u_M \tag{Whiskering}\\
       &= (F_cM(\alpha)(M(t_1'),\ldots, M(t_n'))\bullet \eta_{t_1})u_M\tag{\ref{eq:1}}\\
       &= F_c\bar{p}_M \bullet \eta_{t_1}u_M\tag{Whiskering + Definition of $\bar{p}_M$}
   \end{align*}
   }
  
  \noindent \textbf{Case 3:} Assume $p = B(t_1',\ldots, q,\ldots, t_n')\colon B(t_1',\ldots, t_i',\ldots, t_n')\to B(t_1',\ldots, t_i'',\ldots, t_n')$, where $B\colon \bar{c} = c_1\cdots c_n\to c$ is a functor symbol in $\Sigma$ and the claim holds for the arrow $q\colon t_i'\to t_i''$. Notice that $\eta_{t_1}$ is definitionally the composite:
  {\small$$
  N(B)(N(t_1'),\ldots, N(t_n'))F_0\xRightarrow{N(B)(\eta_{t_1'},\ldots, \eta_{t_n'})} N(B)F_{\bar{c}}(M(t_1'),\ldots, M(t_n'))\xRightarrow{\eta^B(M(t_1'),\ldots, M(t_n'))} F_cM(t_1)
  $$}
  {\small\noindent
  \begin{align*}
       &\eta_{t_2}u_M\bullet \bar{p}_NF^\sigma\\
       &= \eta^B(M^\sigma(t_1'),\ldots, M^\sigma(t_i''),\ldots,  M^\sigma(t_n'))\bullet N(B)(\eta_{t_1'}u_M,\ldots,\eta_{t_i''}u_M,\ldots, \eta_{t_n'}u_M)\bullet\\
       & N(B)(N^\sigma(t_1')F^\sigma,\ldots,\bar{q}_N F^\sigma,\ldots, N^\sigma(t_n')F^\sigma) \tag{Definitions of $\eta_{t_2}$ and $\bar{p}_N$}\\
       &=  \eta^B(M^\sigma (t_1'),\ldots, M^\sigma(t_n'))\bullet N(B)(\eta_{t_1'}u_M, \ldots, \eta_{t_i''}u_M\bullet\bar{q}_NF^\sigma, \ldots, \eta_{t_n'}u_M)\tag{Whiskering}\\
       &= \eta^B(M^\sigma (t_1'),\ldots, M^\sigma(t_n'))\bullet N(B)(\eta_{t_1'}u_M,\ldots, F_{c_i}\bar{q}_M\bullet \eta_{t_i'}u_M,\ldots, \eta_{t_n'}u_M)\tag{Induction on $q$}\\
       &= \eta^B(M^\sigma (t_1'),\ldots, M^\sigma(t_n'))\bullet N(B)F_{\bar{c}}(M^\sigma(t_1'),\ldots,\bar{q}_M,\ldots, M^\sigma(t_n'))\bullet \\
       & N(B)(\eta_{t_1'}u_M,\ldots, \eta_{t_n'}u_M)\tag{Vertical composition componentwise}\\
       &= F_cM(B)(M^\sigma(t_1'),\ldots, \bar{q}_M,\ldots, M^\sigma(t_n'))\bullet \eta^B(M^\sigma(t_1'),\ldots, M^\sigma(t_n'))\bullet\\
       &N(B)(\eta_{t_1'}u_M,\ldots, \eta_{t_n'}u_M)\tag{Interchange}\\
       &= F_c\bar{p}_M\bullet \eta_{t_1}u_M \tag{Definitions of $\bar{p}_M$ and $\eta_{t_1}$}
  \end{align*}
  }

  Thus, the equation \eqref{eq:3} holds for all arrows in $G_\sigma$, and it remains to show the equation for paths. Assume $p$ is a path $p_n\cdots p_1$ of arrows in $G_\sigma$. We show the case $n = 2$, since the general claim follows by an easy induction on $n$, and we denote $p\colon t_1'\xrightarrow{p_1}t_2'\xrightarrow{p_2}t_3'$. Now
  \begin{align*}
       &\eta_{t_3'}u_M\bullet \bar{p}_NF^\sigma\\
       &= \eta_{t_3'}u_M\bullet \bar{p_2}_NF^\sigma \bullet \bar{p_1}_NF^\sigma \tag{$p = p_2p_1$ + Whiskering}\\
       &= F_c\bar{p_2}_M\bullet \eta_{t_2'}u_M\bullet \bar{p_1}_NF^\sigma \tag{Equation \eqref{eq:3} holds for $p_2$}\\
       &= F_c\bar{p_2}_M\bullet F_c\bar{p_1}_M\bullet \eta_{t_1'}u_M \tag{Equation \eqref{eq:3} holds for $p_1$}\\
       &= F_c\bar{p}_M\bullet \eta_{t_1'}u_M \tag{$p = p_2p_1$ + Whiskering}
  \end{align*}\qedhere

\end{proof}
\begin{theorem}[Induced Functor of Algebras Theorem]\label{thm:induced_morphism_between_algebras}
   Let $(F,\eta)\colon M\to N$ be a lax morphism of $\Sigma$-models in $\K$. Assume the following:
\begin{enumerate}
   \item $\eta_t\colon N(t)F_0\Rightarrow F_cM(t)$ is an isomorphism for all function coarity terms $t\colon c$.
   \item $\eta_t*u_Me_M$ is monic as a morphism $N(t)F_0u_Me_M\Rightarrow F_cM(t)u_Me_M$ for every equational coarity term $t\colon c$.
\end{enumerate}
Then there is a unique morphism $F^T\colon M^T\to N^T$, where $e_N F^T = F^\sigma e_M$.
\end{theorem}
\begin{proof}
   The first assumption with Lemma~\ref{lem:morphism_between_inserters}(1) yields the unique morphism $F^\sigma\colon M^\sigma\to N^\sigma$. Consider the following diagram, where the indexing is over $(p,q\colon t_1\to t_2\colon c)\in T$:
   $$ 
\begin{tikzcd}[ampersand replacement=\&]
 {M^T} \& {M^\sigma} \&\& {\prod_{(p,q)} M_c} \\
 \\
 {N^T} \& {N^\sigma} \&\& {\prod_{(p,q)} N_c}
 \arrow["{{{{e_M}}}}", from=1-1, to=1-2]
 \arrow["{{\exists! F^T}}"{description}, dashed, from=1-1, to=3-1]
 \arrow[""{name=0, anchor=center, inner sep=0}, "{{{{(M^\sigma(t_1))}}}}", shift left=3, from=1-2, to=1-4]
 \arrow[""{name=1, anchor=center, inner sep=0}, "{{{{(M^\sigma(t_2))}}}}"', shift right=3, from=1-2, to=1-4]
 \arrow["{{{F^\sigma}}}"', from=1-2, to=3-2]
 \arrow["{{{\prod_{(p,q)}F_c}}}", from=1-4, to=3-4]
 \arrow["{{e_N}}"', from=3-1, to=3-2]
 \arrow[""{name=2, anchor=center, inner sep=0}, "{{{{(N^\sigma(t_1))}}}}", shift left=3, from=3-2, to=3-4]
 \arrow[""{name=3, anchor=center, inner sep=0}, "{{{{(N^\sigma(t_2))}}}}"', shift right=3, from=3-2, to=3-4]
 \arrow["{{{{(\bar{p}_M)}}}}"', shift right=3, Rightarrow, from=0, to=1]
 \arrow["{{{{(\bar{q}_M)}}}}", shift left=3, Rightarrow, from=0, to=1]
 \arrow["{{{{(\bar{q}_N)}}}}", shift left=3, Rightarrow, from=2, to=3]
 \arrow["{{{{(\bar{p}_N)}}}}"', shift right=3, Rightarrow, from=2, to=3]
\end{tikzcd}
$$
   To induce the morphism $F^T$ using the equifier, we must show that $\bar{p}_N*F^\sigma e_M = \bar{q}_N*F^\sigma e_M$ for every $(p,q)\in T$. Assume $(p,q\colon t_1\to t_2\colon c)\in T$. Now
   \begin{align*}
       \eta_{t_2}u_Me_M\bullet \bar{p}_NF^\sigma e_M
       &= F_c\bar{p}_Me_M\bullet \eta_{t_1}u_Me_M \tag{\eqref{eq:3}}\\
       &= F_c\bar{q}_Me_M\bullet \eta_{t_1}u_Me_M \tag{Equifier property of $e_M$}\\
       &= \eta_{t_2}u_Me_M\bullet \bar{q}_NF^\sigma e_M. \tag{\eqref{eq:3}}
   \end{align*}
   By the second assumption $\eta_{t_2}*u_Me_M$ is monic and hence $\bar{p}_N*F^\sigma e_M = \bar{q}_N*F^\sigma e_M$. Thus the equifier $N^T$ induces the unique morphism $F^T\colon M^T\to N^T$, where $e_NF^T = F^\sigma e_M$. 
\end{proof}
\begin{corollary}
   Let $(F,\eta) \colon M \to N$ be either a lax or oplax morphism of $\Sigma$-models in $\K$. 
   \begin{enumerate}
       \item If $F$ is lax and $T$ is algebraic, then $F$ induces a $1$-cell $F^T \colon M^T \to N^T$ in $\K$.
       \item If $F$ is oplax and $T$ is coalgebraic, then $F$ induces a $1$-cell $F^T \colon M^T \to N^T$ in $\K$.
       \item If $F$ is strong, then $F$ induces a $1$-cell $F^T \colon M^T \to N^T$ in $\K$.
   \end{enumerate}
\end{corollary}
\begin{proof}
The lax case follows from Theorem~\ref{thm:induced_morphism_between_algebras}, since if $T$ is algebraic, then $\eta_t$ is an identity for each coarity term $t$ as coarity terms are sorts. The oplax case follows dually. The strong case follows from Theorem~\ref{thm:induced_morphism_between_algebras}, since we inductively see that $\eta_t$ is an isomorphism for each term $t$. 
\end{proof}

As the theory of monoids is algebraic, we recover the following classical result:
\begin{corollary}
   Let $F\colon M\to N$ be a lax monoidal functor between monoidal categories. Then there is an induced functor $\bar{F}\colon \Mon(M)\to \Mon(N)$ between the categories of internal monoids.
\end{corollary}

\section{Internal Limits}
In this section, we talk about limits within an object $c$ in a $2$-category $\K$. A natural notion for a limit in $c$ is a right Kan extension of a morphism $d\colon I\to c$ along a morphism $k\colon I\to K$. Ordinary limits are attained when $\K$ is the $2$-category $\CAT$ and $K$ is the terminal category. Even weighted limits in the enriched context are captured using the formalism of Kan extensions \cite{Kelly:Enriched}.

Understanding the lifting of Kan extensions along PIE limits allows us to prove completeness properties of $0$-cells $M^T$ of bialgebras. This permits us to prove accessibility and local presentability for categories $M^T$ with suitable assumptions on $M$ and $T$.

\begin{definition}
   Let $\C$ be a category. We call an object $c$ in $\C$ a \textbf{quotient initial} object if for any object $c'$ in $\C$ there exists at most one morphism $c\to c'$. \textbf{Subterminality} of an object is defined dually.
\end{definition}
\begin{definition}[Kan Extensions]
   Let $K\xleftarrow{k}I\xrightarrow{d}c$ be morphisms in a $2$-category $\K$. 
       We define the category $\LE(d, k)$ of \textbf{left extensions} $(l\colon K\to c,\varepsilon\colon d\Rightarrow lk)$ of $d$ along $k$ as the comma category $d\downarrow k^*$, where $k^*$ is the pre-composition functor $\K(K,c)\to \K(I,c)$. An initial object of $\LE(d,k)$ is called a \textbf{left Kan extension }of $d$ along $k$. Dually, we define $\RE(d,k)$ of \textbf{right extensions} of $d$ along $k$ as the comma category $k^*\downarrow d$. We call a terminal object in $\RE(d,k)$ a \textbf{right Kan extension} of $d$ along $k$.
\end{definition}

\begin{definition}
   Let $F = (F_i\colon\C\to \D_i)_{i\in I}$ be a family of functors. We say that $F$
   \begin{enumerate}
       \item \textbf{preserves} initial objects, if $F_ic$ is an initial object of $\D_i$ for each initial object $c$ in $\C$ and every $i\in I$.
       \item \textbf{jointly strictly creates} initial objects, if for any initial object $(d_i)_i$ in $\prod_i \D_i$ there is a unique object $c$ in $\C$ so that $F_ic = d_i$ for all $i\in I$ and moreover, $c$ is an initial object in $\C$. In the case $I = \{*\}$, we identify $F$ with $F_*$ and say that $F_*$ \textbf{strictly creates} initial objects.
   \end{enumerate}
   We define the preservation and (joint) strict creation of the terminal object dually. 
\end{definition}
\begin{definition}[Preservation and Strict Creation of Kan Extensions]
   Consider morphisms $K\xleftarrow{k}I\xrightarrow{d}c\xrightarrow{f}c'$ in a $2$-category $\K$. We say that $f$ preserves or strictly creates the left Kan extension of $d$ along $k$ if the induced functor $f_*\colon \LE(d,k)\to \LE(fd,k), \eta\mapsto f*\eta,$ preserves or strictly creates the initial object, respectively. Dually, we define the preservation and strict creation of right Kan extensions. Joint creation of left/right Kan-extensions of $K\xleftarrow{k}I\xrightarrow{d}c$ by a family $(f_i\colon c\to d_i)_i$ of morphisms is defined similarly.
\end{definition}

The following theorem serves as the primary technical result of this section, establishing how the existence of Kan extensions lifts along PIE limits. This theorem provides a unified $2$-categorical framework for limit and colimit inheritance. Consequently, classical results showing that models of universal algebraic theories\footnote{Universal algebraic theory $T$ refers, from our perspective, to an algebraic theory, where the associated $2$-signature is the cartesian signature $\Sigma_{\text{Cart}}$ in Example \ref{ex:$2$-signatures}.}, Eilenberg-Moore categories, and categories of coalgebras inherit structures from their background categories follow as immediate consequences of these three basic $2$-categorical constructions.

\begin{theorem}[Lifting Kan Extensions Along PIE limits]\label{thm:pie_kan_creation}
   Let $\K$ be a $2$-category. Then the following assertions hold:
   \begin{enumerate}
       \item Let $(c_i)_{i \in J}$ be a family of objects in $\K$ with the product $c = \prod_{i\in J} c_i$ existing in $\K$. Then the family $(\pi_i \colon c \to c_i)_{i\in J}$ of projections preserves and jointly strictly creates all right and left Kan extensions.
       \item Consider the inserter diagram
       $$
\begin{tikzcd}
{\Ins(f,g)} \arrow[d, "u"'] \arrow[r, "u"]        & a \arrow[d, "g"] \\
a \arrow[r, "f"'] \arrow[ru, "\phi" description, Rightarrow] & b               
\end{tikzcd}
$$
       in $\K$. Then the canonical morphism $u$ strictly creates left Kan extensions that $f$ preserves, and right Kan extensions that $g$ preserves.
       \item Let $\Eq(\eta,\theta)\xrightarrow{u}c$ be the equifier of $\eta,\theta \colon f \Rightarrow g \colon c \to c'$. Then $u$ strictly creates left (right) Kan extensions whose quotient initiality (subterminality) is preserved by $f$ $(g)$. 
   \end{enumerate}
\end{theorem}
\begin{proof} We show the claims for left Kan extensions as the cases for right Kan extensions follow dually.
   \begin{enumerate}

   \item Let $(l,\varepsilon\colon d\Rightarrow lk)$ be the left Kan extension of $d$ along $k$. Let $i\in J$. We show that $(\pi_il, \pi_i*\varepsilon)$ is an initial left extension. Let $(t_i,\theta_i)$ be an object in $\LE(\pi_id,k)$. Extend 
       $$
       t_j = \begin{cases}
           t_i,\text{ if } i = j,\\
           \pi_jl, \text{ if } i\neq j
       \end{cases}\quad\text{and}\quad\theta_j =\begin{cases}
           \theta_i\colon \pi_id\Rightarrow t_ik\colon I\to c_i,\text{ if }i = j\\
           \pi_j*\varepsilon\colon \pi_jd\Rightarrow t_jk\colon I\to c_j,\text{ if } i\neq j
       \end{cases}
       $$
       for $j\in J$. The second universal property of the product $\prod_i c_i$ induces the transformation $\theta = (\theta_j)_j\colon d\Rightarrow tk$, where $t = (t_j)_j$. By initiality, we have a unique morphism $\phi\colon (l,\varepsilon)\to (t, \theta)$ in $\LE(d,k)$. Thus $\pi_i*\phi$ is a morphism $(\pi_il, \pi_i*\varepsilon)\to (t_i,\theta_i)$ of left extensions. 

       To show uniqueness, let $\alpha\colon (\pi_il,\pi_i*\varepsilon)\to (t_i,\theta_i)$ be any morphism in $\LE(\pi_id,k)$. We define a family of morphisms $\alpha_j \colon \pi_j l \Rightarrow t_j$ by setting $\alpha_i \coloneqq \alpha$ and $\alpha_j \coloneqq \id_{\pi_j l}$ for $j \neq i$. The universal property of the product yields a unique $2$-cell $\psi \colon l \Rightarrow t$ such that $\pi_j * \psi = \alpha_j$ for all $j \in J$. Since $(\alpha_j * k) \bullet (\pi_j * \varepsilon) = \theta_j$ holds for all components $j \in J$, the uniqueness of $2$-cells into a product ensures that $\psi$ is a morphism of left extensions $(l, \varepsilon) \to (t, \theta)$. By the initiality of $(l,\varepsilon)$, we must have $\psi = \phi$, which implies $\alpha = \pi_i * \phi$. Thus, $\pi_i * \phi$ is the unique morphism, confirming that $(\pi_i l, \pi_i * \varepsilon)$ is initial. Hence, $\pi_i$ preserves left Kan extensions. The jointly strict creation of Kan extensions along the projections follows directly from the universal properties of the product $\prod_i c_i$.
       
       \item Let $(l',\varepsilon')$ be a left Kan extension of $ud$ along $k$ that is preserved by $f$. We show that there is a unique lift $(l,\varepsilon)$ in $\LE(d,k)$ so that $l' = ul$ and $\varepsilon' = u*\varepsilon$ and furthermore, $(l,\varepsilon)$ is the left Kan extension of $d$ along $k$. Assume that $(l_1,\varepsilon_1)$ and $(l_2,\varepsilon_2)$ are lifts of $(l',\varepsilon')$ along $u$. We show that $(l_1,\varepsilon_1) = (l_2,\varepsilon_2)$. Note that $\phi*l_1 = \phi*l_2$ as both of these define morphisms $(fl',f*\varepsilon')\to (gl',(g*\varepsilon')\bullet(\phi*d))$, which is seen by the following for $i = 1,2$:
       \begin{align*}
           (\phi*l_ik)\bullet (f*\varepsilon')
           &= (\phi*l_ik)\bullet (fu*\varepsilon_i)\tag{$\varepsilon' = u*\varepsilon_i$}\\
           &= (gu*\varepsilon_i)\bullet (\phi*d)\tag{Bifunctoriality of $*$}\\
           &= (g*\varepsilon')\bullet(\phi*d)\tag{$\varepsilon' = u*\varepsilon_i$}
       \end{align*}
       Thus  $\phi*l_1 = \phi*l_2\colon fl'\Rightarrow gl'$ by the initiality of $(fl',f*\varepsilon')$ in $\LE(fud,k)$. By the first universal property of the inserter $\Ins(f,g)$ there is a unique $l\colon K\to \Ins(f,g)$ so that $l' = ul$ and $\phi*l_1 = \phi*l_2 = \phi*l$. Thus $l_1 = l_2$ and the faithfulness of $u*(-)$ shows that $\varepsilon_1 = \varepsilon_2$. 
       
       
       Next, we construct the lift $(l,\varepsilon)$ of $(l',\varepsilon')$ along $u$. Using the initiality of $(fl',f*\varepsilon')$ in $\LE(fud,k)$, we define $\eta\colon fl'\Rightarrow gl'$ as the unique morphism 
       \begin{equation}\label{cd:1}
           (fl',f*\varepsilon'\colon fud\Rightarrow fl'k)\xrightarrow{\eta} (gl',fud\xRightarrow{\phi*d}gud\xRightarrow{g*\varepsilon'}gl'k).
       \end{equation}
       The first universal property of the inserter $\Ins(f,g)$ with the transformation $fl'\xRightarrow{\eta} gl'$ induces a unique morphism $l\colon K\to \Ins(f,g)$ so that $\eta = \phi*l$ and $ul = l'$. Consider the transformation $\varepsilon'\colon ud\Rightarrow ulk$. Now
       \begin{align*}
           (\phi* lk)\bullet (f*\varepsilon')
           &= (\eta*k)\bullet(f*\varepsilon') \tag{$\eta = \phi*l$}\\
           &= (g*\varepsilon')\bullet (\phi*d). \tag{$\eta$ is a morphism of left extensions \ref{cd:1}}
       \end{align*}
       Thus, the second universal property of the inserter $\Ins(f,g)$ induces a transformation $\varepsilon\colon d\Rightarrow lk$ so that $\varepsilon' = u*\varepsilon$. Thus $(l,\varepsilon)$ is the unique lift of $(l',\varepsilon')$ along $u$.
       
       Let $(t,\theta)$ be a left extension of $d$ along $k$. Lastly, we show that there exists a unique morphism $\lambda\colon (l,\varepsilon)\to(t,\theta)$ in $\LE(d,k)$. Consider the unique morphism $\lambda'\colon (l',\varepsilon')\to (ut,u*\theta)$ in $\LE(ud, k)$. The uniqueness of $\lambda$ is witnessed by the facts that $u*(-)$ is faithful and $u*\lambda = \lambda'$ as $u*\lambda\colon (l',\varepsilon')\to (ut,u*\theta)$ is a morphism in $\LE(ud,k)$. It is sufficient to show that $\lambda'$ has a lift $\lambda\colon l\Rightarrow t$ along $u$, since then faithfulness of $u*(-)$ shows that $\lambda$ is a morphism $(l,\varepsilon)\to(t,\theta)$. For constructing the lift $\lambda$ of $\lambda'$ along $u$ it suffices to show that  $(g*\lambda')\bullet(\phi*l) = (\phi*t)\bullet(f*\lambda')$ holds and then apply the second universal property of the inserter $\Ins(f,g)$. Since $\lambda'\colon (l',\varepsilon')\to (ut,u*\theta)$ is a morphism in $\LE(ud,k)$ and $*$ is a bifunctor, the diagram
       $$
\begin{tikzcd}
                                           & fud \arrow[ld, "f*\varepsilon'"', Rightarrow] \arrow[d, "fu*\theta" description, Rightarrow] \arrow[r, "\phi*d" description, Rightarrow] \arrow[rr, "fu*\varepsilon = f*\varepsilon'", Rightarrow, bend left] & gud \arrow[rd, "gu*\varepsilon" description, Rightarrow] \arrow[d, "gu*\theta" description, Rightarrow] & fulk \arrow[d, "\phi*lk", Rightarrow]      \\
fulk \arrow[r, "f*\lambda'*k"', Rightarrow] & futk \arrow[r, "\phi*tk"', Rightarrow]                                                                                                                                                       & gutk                                                                                                    & gulk \arrow[l, "g*\lambda'*k", Rightarrow]
\end{tikzcd}        $$
       commutes. Thus, the transformations $(\phi*t)\bullet(f*\lambda')$ and $(g*\lambda')\bullet(\phi*l) $ define morphisms
       $$
       (fl',f*\varepsilon')\to (gut, fud\xRightarrow{\phi*d}gud\xRightarrow{gu\theta}gutk)
       $$
       and hence $(g*\lambda')\bullet(\phi*l) = (\phi*t)\bullet(f*\lambda')$ by the initiality of $(fl',f*\varepsilon')$ in $\LE(fud,k)$.

       \item Let $u\colon \Eq(\eta,\theta)\to c$ be the equifier of $\eta,\theta\colon f\Rightarrow g\colon c\to c'$. Let $K\xleftarrow{k}I\xrightarrow{d} \Eq(\eta,\theta)$ be a span in $\K$. Assume that $(l',\varepsilon')$ and $(fl',f*\varepsilon')$ form an initial object in $\LE(ud,k)$ and a quotient initial object in $\LE(fud,k)$, respectively. We show $u$ strictly creates the left Kan extension of $d$ along $k$. The first universal property of equifier shows that $l'$ has at most one lift along $u$, and the second then shows similarly for $\varepsilon'$. To construct the lift $l$ of $l'$ along $u$, it suffices to show that $\eta*l' = \theta*l'$. As $\eta*u = \theta *u$, we have by bifunctoriality of the horizontal composition that the two squares in the diagram
       $$
       \begin{tikzcd}[column sep = 2 cm]
fud \arrow[d, "\eta*ud = \theta*ud" description, Rightarrow] \arrow[r, "f*\varepsilon'", Rightarrow] & fl'k \arrow[d, "\theta*l'k", shift left, Rightarrow] \arrow[d, "\eta*l'k"', shift right, Rightarrow] \\
gud \arrow[r, "g*\varepsilon'"', Rightarrow]                                              & gl'k                                                                         
\end{tikzcd}
       $$
       commute. This shows that $\eta*l'$ and $\theta*l'$ define morphisms $(fl',f*\varepsilon')\to (gl', fud\xRightarrow{\eta*ud}gud\xRightarrow{g\varepsilon'}gl'k)$ in $\LE(fud,k)$ and thus by the quotient initiality of $(fl',f*\varepsilon')$ we have that $\eta*l = \theta*l$. Hence $(l',\varepsilon')$ has a lift $(l,\varepsilon)$ along $u$. The full faithfulness of $u*(-)$ shows that $(l,\varepsilon)$ is initial in $\LE(d,k)$. \qedhere

   \end{enumerate}
\end{proof}
\begin{definition}
Given a signature pair $(\Sigma,\sigma = (S,F))$ a $\sigma$-theory $T$ and a $\Sigma$-model $M$ in a PIE complete $2$-category, we call the morphism $M(t)\colon M_0\to M_c$ a \textbf{function/equational arity/coarity support}, if $t\colon c\in \Term_S^\Sigma$ is a function/equational arity/coarity term, respectively, relative to $T$.
\end{definition}
\begin{corollary}\label{cor:LimitsLiftedToBialgebras}
   Let $(\Sigma,\sigma)$ be a signature pair with a $\sigma$-theory $T$. Let $M$ be a $\Sigma$-model within a PIE-complete $2$-category $\K$. Consider the canonical morphism $u\colon M^T\to M_0$. Then the following assertions hold:
   \begin{enumerate}
       \item The morphism $u$ strictly creates those left Kan extensions that function supports $M(t_1)$ preserve and whose quotient initiality equational supports $M(t_1')$ preserve.
       \item The morphism $u$ strictly creates those right Kan extensions that function coarity supports $M(t_2)$ preserve and whose subterminality equational coarity supports $M(t'_2)$ preserve. 
       \item If $T$ is algebraic, then $u$ strictly creates all right Kan extensions.
       \item If $T$ is coalgebraic, then $u$ strictly creates all left Kan extensions.
   \end{enumerate}
\end{corollary}
\begin{proof}
   It suffices to only show the first case, as the second follows dually, and the third and fourth follow from the fact that projections preserve all left and right Kan extensions by Theorem~\ref{thm:pie_kan_creation}(1).

   Note that $M^\sigma = \Ins(L_1,L_2)$, where $L_i= (M(t_i))_{f\colon t_1\to t_2}\colon M_0\to \prod_f M_{c_f}$. The morphism $L_1$ preserves exactly those left Kan-extensions that morphisms $M(t_1)$, for $\sigma$-function symbols $f\colon t_1\to t_2$, preserve by Theorem~\ref{thm:pie_kan_creation} (1). Thus $u_M\colon M^\sigma\to M_0$ strictly creates those left Kan extensions that $M(t_1)$ preserves for each $\sigma$-function symbol $f\colon t_1\to t_2$ by Theorem $\ref{thm:pie_kan_creation}(2)$. Similarly, using Theorem~\ref{thm:pie_kan_creation}(3), we have that the canonical morphism $M^T\to M^\sigma$ strictly creates those left Kan extensions whose quotient initiality is preserved by $M^\sigma(t_1')\colon M^\sigma\to M_0\xrightarrow{M(t_1')} M_{c'}$ for every equation $(p,q\colon t_1'\to t_2')\colon c'\in T$. Whence, the claim follows.
\end{proof}
\begin{corollary}\label{cor:limitPreservationFromAtomicPreservation}
   Let $(\Sigma,\sigma)$ be a signature pair with a $\Sigma$-model $M$ in a PIE-complete $2$-category $\K$ and a $\sigma$-theory $T$. Let $k\colon I\to K$ be a morphism in $\K$. Assume $M(B)$ preserves all left Kan extensions along $k$ for each functor symbol $B$ in $\Sigma$. Let $t\colon c$ be a term. Then $M(t)\colon M_0\to M_c$ preserves all left Kan extensions along $k$. In particular, the canonical morphism $u\colon M^T\to M_0$ strictly creates all left Kan extensions along $k$.
\end{corollary}
\begin{proof}
   We proceed by structural induction on the term $t$. Let $d\colon I\to M_0$ be a morphism and assume the left Kan extension of $d$ along $k$ exists. If $t$ is a sort $s$, the term functor is a canonical projection $M_0 \to M_{c_s}$. Projections preserve all Kan extensions by Theorem~\ref{thm:pie_kan_creation}(1), so the base case holds.

   Assume then that $t = B(t_1,\dots, t_n)$ and that $M(t_i)$ preserves left Kan extensions along $k$ for each $i\leq n$. Now
   $$
   M(t) = M(B)\circ (M(t_1),\dots, M(t_n))
   $$
   and since by assumption $M(t_i)$ preserves left Kan extensions along $k$, it follows from Theorem~\ref{thm:pie_kan_creation}(1) that $(M(t_1),\dots, M(t_n))$ preserves them as well. By assumption, $M(B)$ preserves left Kan extensions along $k$ and therefore the composite $M(t)$ preserves them also.

   The last part follows from Corollary~\ref{cor:LimitsLiftedToBialgebras}(1).
\end{proof}

Consider a sifted category $I$, a non-empty category where the diagonal $\Delta\colon I\to I\times I$ is final \cite{AdamekRosickyVitale2010SiftedColimits}.\footnote{Recall that a functor $F\colon \C\to \D$ is final if $d\downarrow F$ is a connected (implying non-empty) category for each object $d$ in $\D$ and equivalently for any diagram $H\colon \D\to \E$ the functor of cocones $F^*\colon \mathrm{Cocone}(H)\to \mathrm{Cocone}(HF)$ is an isomorphism of categories.} All directed posets and the category for a reflexive pair $\begin{tikzcd}
\bullet \arrow[r, shift left=2] \arrow[r, shift right=2] & \bullet \arrow[l]
\end{tikzcd}$, with the middle arrow defining a joint section, form sifted categories.

\begin{lemma}\label{lem:SiftedTotalityFromComponents}
   Let $F\colon \C_1\times\ldots\times \C_n\to \C$ be a functor and let $I$ be a sifted category. Then $F$ preserves $I$-colimits if and only if $F$ preserves $I$-colimits componentwise, meaning the functor $x_i\mapsto F(c_1,\ldots, x_i,\ldots,c_n)\colon \C_i\to \C$ preserves $I$-colimits for  $c_j$ in $\C_j$ and $i,j\leq n$. 
\end{lemma}
\begin{proof}
   Clearly, the claim holds for $n = 0,1$. The constant functors $1\to \C_i$ preserve connected colimits and, therefore, so do the functors $\C_i\to \C_1\times \ldots \times \C_n,x\mapsto (c_1,\ldots, x,\ldots, c_n)$. Thus, if $F$ preserves $I$-colimits, then so do the functors $F\circ (c_1,\ldots, -,\ldots, c_n)$. Let us then assume that $F$ preserves $I$-colimits componentwise. We show that $F$ preserves $I$-indexed colimits. It suffices to prove the claim assuming $n = 2$, since the induction step is then easy. Consider $D = (D_1,D_2)\colon I\to \C_1\times \C_2$ with a colimit. Now 
   \begin{align*}
       F(\colim_i (D_1(i), D_2(i)))
       &\cong F(\colim_{i,j}(D_1(i),D_2(j)))\tag{$\Delta$ is final}\\
       &\cong F(\colim_i \colim_j(D_1(i),D_2(j)))\tag{Fubini}\\
       &\cong \colim_i F(\colim_j(D_1(i),D_2(j)))\tag{Preservation componentwise}\\
       &\cong \colim_i\colim_j F(D_1(i),D_2(j))\tag{Preservation componentwise}\\
       &\cong \colim_{i,j} F(D_1(i),D_2(j))\tag{Fubini}\\
       &\cong \colim_i F(D_1(i),D_2(i)).\tag{$\Delta$ is final}
   \end{align*}
\end{proof}

The following corollary highlights the ubiquity of sifted colimits and cosifted limits in practice. Whenever a monoidal category possesses reflexive coequalizers that are preserved componentwise by its tensor product, any category of internal bialgebras constructed over it strictly inherits these reflexive coequalizers.

\begin{corollary}\label{cor:liftingSifted}
   Let $(\Sigma, \sigma)$ be a signature pair with $\sigma$-theory $T$. Let $M$ be a categorical $\Sigma$-model and $I$ be a sifted category. If $M(B)$ preserves $I$-colimits componentwise for each functor symbol $B$ in $\Sigma$, then the forgetful functor $U\colon M^T\to M_0$ strictly creates all $I$-indexed colimits. 
\end{corollary}
\begin{proof}
   Follows immediately from Lemma~\ref{lem:SiftedTotalityFromComponents} and Corollary~\ref{cor:limitPreservationFromAtomicPreservation}.
\end{proof}
In the following, we list some famous results that follow from our framework.
\begin{corollary}
   Let $T$ be a monad on a category $\C$. Then the forgetful functor $\C^T\to \C$ from the Eilenberg-Moore category $\C^T$ of $T$ strictly creates limits and those colimits that $T$ preserves and whose quotient initiality $T^2$ preserves.
\end{corollary}

\begin{corollary}
   Let $F\colon \C\to \E$ and $G\colon \D\to \E$ be functors. Then the forgetful functor $U\colon F\downarrow G\to \C\times \D$ strictly creates those colimits and limits that $F$ and $G$ preserve, respectively.
\end{corollary}

\begin{corollary}
   Let $c$ be an object of a category $\C$. Then the forgetful functor $U\colon \C/c\to \C$ from the slice category over $c$ strictly creates all colimits and connected limits.
\end{corollary}
\begin{proof}
   The theory for $C/c$ is coalgebraic, which yields the creation of colimits. Since the constant functor $c\colon 1\to C$ preserves all connected limits, $\C/c\to \C$ strictly creates connected limits. 
\end{proof}
\begin{corollary}
    The category $\mathbf{Field}$ of fields has connected limits and sifted colimits.
\end{corollary}
\begin{proof}
    Consider the signature pair $(\Sigma,\sigma)$ with the theory of fields $T$ and the categorical $\Sigma$-model $M = (\Set, \times, 1, (-)+1)$ from Example~\ref{exs:theories}(3). Note that $M(B)$ preserves connected limits (the product $\times\colon \Set^2\to \Set$ and the constant $1$ are right adjoints, and $x\mapsto x+1\colon \Set\to \Set$ preserves connected limits) and sifted colimits (the product is a left adjoint componentwise and hence preserves sifted colimits by Lemma~\ref{lem:SiftedTotalityFromComponents}, the constant preserves connected colimits, and the functor $(-)+1$ preserves all colimits) for each $\Sigma$-functor symbol $B$. Thus, the equivalence $M^T\simeq \mathbf{Field}$ together with Corollary~\ref{cor:limitPreservationFromAtomicPreservation} establishes the claim.
\end{proof}

\subsection{Accessibility and Local Presentability of Bialgebras}
Let us fix a signature pair $(\Sigma,\sigma)$, a $\sigma$-theory $T$, and a locally presentable symmetric monoidal closed category $\V$. We denote by $\ACC_\V$ the $2$-category of large $\V$-accessible categories, $\V$-accessible functors and $\V$-natural transformations \cite{LACK2023107196}. We establish one of our main results, showing how accessibility and local presentability lift to the object of $T$-bialgebras. By working directly in the enriched setting, we obtain the ordinary categorical results as a simple corollary.

\begin{theorem}[$\V$-Enriched Accessibility of $M^T$]\label{thm:enriched_accessability_lifting}
    Let $M$ be a $\Sigma$-model in the $2$-category $\CAT_{\V}$ of $\V$-enriched categories. If each $\V$-category $M_c$ is $\V$-accessible and each $\V$-functor $M(B)$ is $\V$-accessible, then the $\V$-category $M^T$ of $T$-bialgebras is $\V$-accessible and the forgetful $\V$-functor $U\colon M^T\to M_0$ is $\V$-accessible.
\end{theorem}
\begin{proof}
   The $2$-category $\mathbf{ACC}_{\V}$ of $\V$-accessible categories, $\V$-accessible functors, and $\V$-natural transformations is closed under strict PIE limits in $\CAT_{\V}$ by Theorem~5.5 in \cite{LACK2023107196}. Thus, $U\colon M^T\to M_0$ is $\V$-accessible.
\end{proof}


\begin{remark}
   If $\V = \Set$, the functors $M(B)$ in Theorem~\ref{thm:enriched_accessability_lifting} may be equivalently assumed to be accessible componentwise. Componentwise accessibility allows one to extract a uniform regular cardinal $\lambda$ so that $M(B)$ is $\lambda$-accessible in each component for each $\Sigma$-functor symbol $B$ \cite[Example 2.13(6)]{LPC}. As directed colimits are sifted colimits, componentwise preservation of $\lambda$-directed colimits is equivalent to the preservation of $\lambda$-directed colimits by Lemma~\ref{lem:SiftedTotalityFromComponents}.
\end{remark}

\begin{theorem}[Local Presentability of $M^T$]\label{thm:LiftingLP}
   Let $M$ be a $\Sigma$-model in the $2$-category $\CAT_{\V}$ where each $\V$-category $M_c$ is $\V$-locally presentable and the functor $M(B)$ is $\V$-accessible for each functor symbol $B$. Consider the $\V$-functor $U\colon M^T\to M_0$. Then the following conditions hold:
   \begin{enumerate}
       \item If each arity support $M(t) \colon M_0 \to M_c$ is cocontinuous, then $M^T\to M_0$ is comonadic and $M^T$ is locally presentable. In particular, if $T$ is a coalgebraic theory, then $M^T$ is locally presentable, and $U$ is comonadic.
       \item If each coarity support $M(t) \colon M_0 \to M_c$ is continuous, then $M^T\to M_0$ is monadic and $M^T$ is locally presentable. In particular, if $T$ is an algebraic theory, then $M^T$ is locally presentable and $U$ is monadic.
   \end{enumerate}
\end{theorem}
\begin{proof}
   We prove the monadic case, as the comonadic case follows similarly. Assume each coarity support $M(t) \colon M_0 \to M_c$ is continuous. Since $M^T$ is $\V$-accessible (Theorem~\ref{thm:enriched_accessability_lifting}) and $\V$-complete (Corollary~\ref{cor:LimitsLiftedToBialgebras}(2)), $M^T$ is $\V$-locally presentable by \cite[Corollary 3.21.]{LACK2023107196}. As the functor $U\colon M^T\to M_0$ is accessible (Theorem~\ref{thm:enriched_accessability_lifting}) and continuous between $\V$-locally presentable categories, the Enriched Adjoint Functor Theorem \cite[Theorem 5.32]{Kelly:Enriched} shows that $U$ has a left adjoint. Lastly, as $U$ satisfies the enriched Beck's monadicity criterion (\cite[Theorem II.2.1]{dubuc1970}), $U$ is $\V$-monadic. 
\end{proof}

If $\C$ is a locally presentable symmetric monoidal category, where the tensor is accessible componentwise, then the category of internal bimonoids forms a locally presentable category. This is seen by the fact that the category of bimonoids is the category $\textbf{Comon}(\Mon(\C))$ and that $\Mon(\C)$ is itself a locally presentable (Theorem~\ref{thm:LiftingLP}) symmetric monoidal category with the original tensor, which also preserves large enough directed colimits componentwise (as these are computed in $\C$).

\section{Factorization Systems, Regularity, and Exactness}

To establish how regularity and exactness lift along the construction $M\mapsto M^T$ in $\CAT$, we first study how orthogonal factorization systems lift along $2$-dimensional limits. We then generalize these notions to $2$-categories via the representable $2$-functors $\K(-,c)\colon \K^{op}\to \CAT$.

\subsection{Orthogonal Factorization Systems}
\begin{definition}[Orthogonal Factorization System]\label{def:1cat_ofs}
   Let $\C$ be a category. Two morphisms $e$ and $m$ in $\C$ are said to be \textbf{orthogonal}, denoted $e \perp m$, if for every commutative square 
   $$
\begin{tikzcd}
a \arrow[d, "e"'] \arrow[r, "f"]                              & c \arrow[d, "m"] \\
b \arrow[r, "g"'] \arrow[ru, "\exists!h" description, dotted] & d               
\end{tikzcd}
   $$ there exists a unique morphism $h\colon c\to d$, a \textbf{diagonal filler}, making the diagram commute. We extend the notation $A\perp B$ to include classes of morphisms $A$ and $B$ in $\C$. We denote $A^\perp\coloneqq \{m\in \Mor(\C)\mid a\perp m\text{ for all $a\in A$}\}$ and similarly for $^\perp\! A$ and call the classes the \textbf{right} and the \textbf{left orthogonal classes} of $A$, respectively.
   
   An \textbf{orthogonal factorization system (OFS)} on $\C$ is a pair $(E, M)$ of classes of morphisms in $\C$ such that:
   \begin{enumerate}
       \item \textbf{Factorization:} Every morphism $f$ in $\C$ admits a factorization $f = m \circ e$ with $e \in E$ and $m \in M$. Such a factorization is called an $(E,M)$-factorization.
       \item \textbf{Orthogonality:} $M = E^\perp$ and $E =\! ^\perp\! M$
   \end{enumerate}
   An OFS $(E,M)$ is called \textbf{stable} if every pullback of any morphism in $E$ is in $E$. 
\end{definition}

\begin{remark}
   An orthogonal factorization system can be equivalently described as a pair $(E, M)$ of classes of morphisms of $\C$ such that every morphism in $\C$ admits an $(E, M)$-factorization, $E \perp M$, and both $E$ and $M$ contain all isomorphisms and are closed under composition and retracts in the arrow category $\C\downarrow \C$ \cite[Lemma 11.2.3]{Riehl2014CategoricalHomotopyTheory}. In practice, these closure properties are useful when verifying that a lifted pair of classes forms an orthogonal factorization system. 
\end{remark}


We now show how these factorization systems strictly lift along PIE limits in $\CAT$.

\begin{theorem}[Lifting OFS Along PIE Limits]\label{thm:pie_ofs_creation}
   Let $\C_i,\C,\D$ be categories with orthogonal factorization systems $(E'_i,M'_i), (E',M'), (E'',M'')$, respectively, for each $i\in I$. Consider functors $F,G\colon \C\to \D$ that preserve the left and right classes, respectively. Then the following hold:
   \begin{enumerate}
       \item $(\prod_iE_i',\prod_i M_i')$ is the unique OFS on $\prod_i \C_i$ preserved by the projections.
       \item $(U^{-1}(E'),U^{-1}(M'))$ is the unique OFS on $\Ins(F,G)$ preserved by the forgetful functor $U\colon \Ins(F,G)\to \C$.
       \item Let $\eta,\theta\colon F\Rightarrow G$ be natural transformations. Then the equifier $\Eq(\eta,\theta)$ has a unique OFS $(U^{-1}(E'), U^{-1}(M'))$ that is preserved by the forgetful functor $U\colon \Eq(\eta,\theta)\to \C$.
   \end{enumerate}
\end{theorem}
\begin{proof}
\hfill
\begin{enumerate}
   \item The product $\prod_i \C_i$ has an orthogonal factorization system defined by the product $(\prod_iE_i',\prod_i M_i')$. If $(E,M)$ is an OFS on $\prod_i \C_i$ preserved by the projections, we would have that $E\subset \prod_i E_i'$ and $M\subset \prod_i M_i'$ and by orthogonality then $\prod M_i'= M$ and hence $E = \prod_i E_i'$ proving the uniqueness.

   \item Define $E \coloneqq U^{-1}(E')$ and $M \coloneqq U^{-1}(M')$. Note that if $(E,M)$ is an orthogonal factorization system on $\Ins(F,G)$, then its uniqueness is seen as in the previous part. Clearly, both $E$ and $M$ are closed under composition, isomorphisms, and retractions; and so it suffices to show that each morphism in $\Ins(F,G)$ has an $(E,M)$-factorization and $E\perp M$.
   
   Let $f\colon (x,\phi_x\colon Fx\to Gx)\to (z,\phi_z\colon Fz\to Gz)$ be a morphism in $\Ins(F,G)$. Now $f$ factors $f\colon x\xrightarrow{e}y\xrightarrow{m}z$ as a morphism in $\C$ for some $e\in E'$ and $m\in M'$. Using the orthogonality $F(e)\perp G(m)$, we define the morphism $\phi_y\colon Fy\to Gy$ as the unique morphism making the diagram
   $$
   \begin{tikzcd}[column sep = 2 cm]
F(x) \arrow[d, "\phi_x"'] \arrow[r, "F(e)\in E''"] & F(y) \arrow[d, "\exists!\phi_y" description] \arrow[r, "F(m)"] & F(z) \arrow[d, "\phi_z"] \\
G(x) \arrow[r, "G(e)"']                              & G(y) \arrow[r, "G(m)\in M''"']                                   & G(z)                      
\end{tikzcd}
   $$
   commute. Thus we have a factorization $f = me$ in $\Ins(f,g)$, where $m\in M$ and $e\in E$. Finally, we show $E\perp M$. Consider the following lifting problem:
   $$
   \begin{tikzcd}
(x,\phi_x) \arrow[d, "e\in E"'] \arrow[r, "f"] & (y,\phi_y) \arrow[d, "m\in M"] \\
(v,\phi_v) \arrow[r, "g"']                      & (w,\phi_w)                     
\end{tikzcd}
   $$
   in $\Ins(F,G)$. Note that there is a unique solution $\theta\colon v\to y$ in the underlying category $\C$. This shows the lifting problem has at most one solution. It suffices to show that $\theta$ is a morphism in $\Ins(F,G)$. We show that the diagram
   $$
   \begin{tikzcd}
F(v) \arrow[d, "\phi_v" description] \arrow[r, "F(\theta)"] & F(y) \arrow[d, "\phi_y" description] \\
G(v) \arrow[r, "G(\theta)"']                                  & G(y)                                  
\end{tikzcd}
   $$
   commutes. The orthogonality $F(e)\perp G(m)$ induces a unique morphism $\phi\colon Fv\to Gy$ making the two parts of the whole exterior of the rectangle in the diagram
   $$
\begin{tikzcd}[column sep = 2 cm]
Fx \arrow[d, "\phi_x" description] \arrow[r, "F(e)" description] \arrow[rr, "Ff", bend left] & Fv \arrow[r, "F(\theta)" description] \arrow[d, "\phi_v" description] \arrow[rd, "\exists!\phi" description] \arrow[rr, "F(g)", bend left] & Fy \arrow[r, "F(m)" description] \arrow[d, "\phi_y" description] & Fw \arrow[d, "\phi_w" description] \\
Gx \arrow[r, "G(e)" description] \arrow[rr, "G(f)"', bend right]                                & Gv \arrow[r, "G(\theta)" description] \arrow[rr, "G(g)"', bend right]                                                                         & Gy \arrow[r, "G(m)" description]                                    & Gw                                   
\end{tikzcd}
$$
   commute. As both morphisms $G(\theta)\phi_v$ and $\phi_yF(\theta)$ satisfy the defining condition of $\phi$, it follows that $G(\theta)\phi_v = \phi_y F(\theta)$. Thus $\theta\colon (v,\phi_v)\to (y,\phi_y)$ is a morphism in $\Ins(F,G)$. 
   
   \item Define $E\coloneqq U^{-1}(E')$ and $M = U^{-1}(M')$. If $(E,M)$ is an orthogonal factorization system, then it is clearly the unique one that is preserved by $U$. We show that $(E,M)$ forms an orthogonal factorization system. Let $f\colon x\to z$ be a morphism in $\Eq(\eta,\theta)$ and consider a factorization $f\colon x\xrightarrow{e}y\xrightarrow{m}z$, where $e\in E'$ and $m\in M'$. By the full faithfulness of $U$, it is sufficient to check that $\eta_y = \theta_y$ to show that $(E,M)$ is an orthogonal factorization system on $\Eq(\eta,\theta)$. 
   Consider the following diagram:
   $$
\begin{tikzcd}[column sep = 2cm]
F(x) \arrow[d, "\eta_x = \theta_x" description] \arrow[r, "F(e)\in E''"] & F(y) \arrow[d, "\eta_y"', shift right] \arrow[d, "\theta_y", shift left] \arrow[r, "F(m)"] & F(z) \arrow[d, "\eta_z = \theta_z" description] \\
G(x) \arrow[r, "G(e)"']                                                 & G(y) \arrow[r, "G(m)\in M''"']                                                              & G(z)                                           
\end{tikzcd}
$$
   As $F(e)\perp G(m)$, we have by orthogonality that the connecting morphism is unique. As both of the morphisms $\eta_y$ and $\theta_y$ make both of the squares commute, we have the equality $\eta_y = \theta_y$. This shows that $y$ is an object of $\Eq(\eta,\theta)$. \qedhere
\end{enumerate}
\end{proof}

\begin{definition}
   Let $\C$ be a category. We say that $\C$ is a \textbf{regular} category, if $\C$ has finite limits and a stable OFS $(E,M)$ where $M$ is the class of monomorphisms. A regular category $\C$ is called \textbf{exact}, if each internal equivalence relation $R\rightrightarrows c$ in $\C$ is a kernel pair of some morphism in $\C$. \footnote{A jointly monomorphic pair of morphisms $R\rightrightarrows c$ in a category $\C$ is called an equivalence relation, if the induced subset of $\hom(x,c)\times \hom(x,c)$ is an equivalence relation for each object $x$ in $\C$.}
\end{definition}
In a category $\C$ with pullbacks, we have that $E =\! ^\perp\! M$ and $M = E^\perp$ where $M$ is the class of monomorphisms and $E$ is the class of extremal epimorphisms in $\C$. In a regular category, each extremal epimorphism is a regular epimorphism by Proposition 2.2 (b) in \cite{Kelly1991}. Therefore, in a regular category, the coequalizers of kernel pairs exist.

\begin{corollary}\label{cor:liftingregularexact}
   Let $\C_i, \C,\D$ be categories for $i\in I$. Let $F,G\colon \C\to \D$ be functors preserving regular epis and finite limits, respectively. Then the following holds:
   \begin{enumerate}
       \item If $\C_i$ is a (exact) regular category for each $i\in I$, then so is $\prod_{i\in I} \C_i$.
       \item If $\C$ and $\D$ are regular categories, then so is $\Ins(F,G)$. Furthermore, if $F$ preserves the coequalizers of kernel pairs, and $\C$ and $\D$ are exact, then $\Ins(F,G)$ is exact as well.
       \item If $\C$ and $\D$ are (exact) regular categories and $\eta,\theta\colon F\Rightarrow G\colon \C\to \D$ are parallel natural transformations, then the equifier category $\Eq(\eta,\theta)$ is (exact) regular.
   \end{enumerate}
\end{corollary}
\begin{proof}
       We only prove the second assertion, as the others are shown similarly. 
       Assume that $\C$ and $\D$ are regular categories. Since $\C$ has, and $G$ preserves, finite limits, it follows that $\Ins(F,G)$ is finitely complete by Theorem~\ref{thm:pie_kan_creation}(2). Consider the classes of morphisms $E = U^{-1}(E')$ and $M = U^{-1}(M')$, where $M'$ and $E'$ are the classes of monomorphisms and extremal epimorphisms in $\C$, respectively. By Theorem~\ref{thm:pie_ofs_creation}(2), $(E,M)$ is an orthogonal factorization system on $\Ins(F,G)$ as $F$ preserves regular epis and $G$ monics. Note that $M$ is the class of monomorphisms of $\Ins(F,G)$ and $E$ is stable as the forgetful functor $U\colon \Ins(F,G)\to \C$ strictly creates limits. Thus $\Ins(F,G)$ is a regular category.

       Assume then that $\C$ and $\D$ are exact and $F$ preserves the coequalizers of kernel pairs. We show that $\Ins(F,G)$ is an exact category. Consider an equivalence relation $R\rightrightarrows x$ in $\Ins(F,G)$. As $U$ preserves limits, it follows that $U(R)\rightrightarrows U(x)$ is an equivalence relation in $\C$. By exactness, we may consider the coequalizer $f\colon Ux\to y$ of the kernel pair $UR\rightrightarrows Ux$. As $F$ preserves the coequalizers of kernel pairs, $F$ preserves the coequalizer of $UR\rightrightarrows Ux$. Thus $U$ strictly creates both the coequalizer of $R\rightrightarrows x$ in $\Ins(F,G)$ and its kernel pair, as $G$ preserves pullbacks. Thus, the equivalence relation $R\rightrightarrows x$ is a kernel pair of $f$ in $\Ins(F,G)$, proving that $\Ins(F,G)$ is exact.
\end{proof}



\subsection{Representable Regularity and Exactness}

One can formulate an orthogonal factorization system on a category $\C$ as a suitable idempotent comonad-monad pair on the arrow category $\C\downarrow\C$, so-called algebraic factorization system \cite{GrandisTholen2006NaturalWFS}. This allows a $2$-categorical formulation of an orthogonal factorization system on an object $c$ in a $2$-category $\K$ with arrow objects. However, as the representable $2$-functors preserve and jointly create the idempotent monad-comonads on arrow objects, we can instead define the notion of an orthogonal factorization system on $c$ representably.

\begin{definition}\label{def:representable_ofs}
   Let $c$ be an object of a $2$-category $\K$. Let $(E_x,M_x)$ be an OFS on the hom-category $\K(x,c)$ for each object $x$ in $\K$. We call the pair $(E,M) = ((E_x)_x,(M_x)_x)$ a \textbf{representable orthogonal factorization system} on $c$, if for every morphism $f\colon x\to y$ in $\K$, the pre-composition functor $f^* \colon \K(y,c) \to \K(x,c), \eta\mapsto \eta*f,$ preserves the left and the right classes.
\end{definition}
We can similarly define the exactness properties of an object representably.

\begin{definition}\label{def:representable_exactness}
   An object $c$ in a $2$-category $\K$ is \textbf{representably regular} if for every object $x$ in $\K$, the hom-category $\K(x, c)$ is a regular category, and for every morphism $f \colon x \to y$, the precomposition functor $f^* \colon \K(y,c) \to\K(x, c)$ preserves finite limits and regular epimorphisms. Furthermore, $c$ is \textbf{representably exact} if it is representably regular and every hom-category of the form $\K(x, c)$ is exact.
\end{definition}
In general, we say that a morphism $f\colon c\to d$ in a $2$-category $\K$ preserves structures or properties (for instance regularity, exactness, factorization systems) \textbf{representably}, if the post composition functor $f_*\colon \K(x,c)\to \K(x,d)$ preserves those properties for each object $x$ in $\K$.

\begin{corollary}\label{cor:pie_ofs_creation_2cat}
       Let $c_i,c,d$ be objects in a $2$-category $\K$ with representable orthogonal factorization systems $(E'_i,M'_i), (E',M'), (E'',M'')$, respectively, for each $i\in I$. Consider morphisms $f,g\colon c\to d$ that representably preserve the left and right classes, respectively. Then the following holds.
   \begin{enumerate}
       \item $(\prod_iE_i',\prod_i M_i') $ is the unique OFS on $\prod_i c_i$, assuming the product exists,  which is preserved by the projections.
       \item Assume that the inserter $\Ins(f,g)$ exists in $\K$. Then $(u_*^{-1}(E'),u_*^{-1}(M'))$ is the unique OFS on $\Ins(f,g)$ preserved by the canonical morphism $u\colon \Ins(f,g)\to c$.
       \item Assume that the equifier $\Eq(\eta,\theta)$ exists. Then it has a unique OFS $(u_*^{-1}(E'), u_*^{-1}(M'))$ that is preserved by the forgetful canonical morphism $u\colon \Eq(\eta,\theta)\to c$.
   \end{enumerate}
\end{corollary}
\begin{proof}
   We only show the first part, as the rest are shown similarly. Assume that the product $\prod_{i\in I} c_i$ exists in $\K$. Let $x$ be an object of $\K$ and note that the representable functor $\K(x,-)\colon \K\to \CAT$ preserves products. Thus we have the projection $\K(x,\pi_i)\colon \K(x,\prod_i c_i)\to \K(x,c_i)$ for $i\in I$. Theorem~\ref{thm:pie_ofs_creation}(1) shows that $\K(x,\prod_i c_i)$ attains the unique orthogonal factorization system as desired. Consider then a morphism $f\colon x\to y$ in $\K$ and the commutative diagram
   $$
   \begin{tikzcd}
{\K(y,\prod_i c_i)} \arrow[d, "f^*"'] \arrow[r, "\cong"] & {\prod_i \K(y,c_i)} \arrow[d, "{\prod_i f^*}"] \\
{\K(x,\prod_i c_i)} \arrow[r, "\cong"']                  & {\prod_i \K(x,c_i)}                                 
\end{tikzcd}
   $$
   As all but the functor $f^*\colon \K(y,\prod_i c_i)\to \K(x,\prod_i c_i)$ are known to preserve the orthogonal factorization systems in the diagram above, it follows that $f_*$ also respects the orthogonal factorization system.
\end{proof}
\begin{corollary}\label{cor:pie_exactness_creation_2cat}
   Let $c_i, c,d$ be objects in a $2$-category $\K$ with PIE limits for $i\in I$. Assume that morphisms $f,g\colon c\to d$ representably preserve regular epis and finite limits, respectively. Then the following holds:
   \begin{enumerate}
       \item If $c_i$ is representably (exact) regular for each $i\in I$, then so is $\prod_{i\in I} c_i$.
       \item If $c$ and $d$ are representably regular objects, then so is $\Ins(f,g)$. Furthermore, if $f$ representably preserves the coequalizers of kernel pairs, and $c$ and $d$ are representably exact, then $\Ins(f,g)$ is representably exact as well.
       \item If $c$ and $d$ are representably (exact) regular and $\eta,\theta\colon F\Rightarrow G\colon \C\to \D$ are parallel transformations, then the equifier $\Eq(\eta,\theta)$ is representably (exact) regular.
   \end{enumerate}
\end{corollary}

\begin{proof}
   The claims follow from the fact that representable $2$-functors $\K(x,-)$ preserve PIE limits and from Corollary~\ref{cor:liftingregularexact}.
\end{proof}

Proposition 3.7 in \cite{Wissmann2022Minimality} shows how a monad $T\colon \C\to \C$ preserving the left class of an orthogonal factorization system $(E,M)$ on $\C$ induces an orthogonal factorization system on the Eilenberg-Moore category $\C^T$, generalizing the original result by Linton in \cite{Linton1969}. We further generalize these results by lifting orthogonal factorization systems to arbitrary categories of bialgebras in the following corollary:

\begin{corollary}\label{cor:StructuresLiftedToBialgebras}
   Let $(\Sigma,\sigma)$ be a signature pair with a $\sigma$-theory $T$. Let $M$ be a $\Sigma$-model within a PIE-complete $2$-category $\K$. Consider the forgetful morphism $u\colon M^T\to M_0$. Then the following assertions hold:
   \begin{enumerate}
       \item Assume $M_c$ is equipped with a representable OFS $(E^c,M^c)$ for each category symbol $c$. If each arity support $M(t_1)$ and coarity support $M(t_2)$ representably preserve the left and right classes, respectively, then there exists a unique representable OFS on $M^T$ representably preserved by the projections $M^T\to M_{c_s}$. In particular, the conclusion follows if $M(B)\colon M_{\bar{c}}\to M_d$ representably preserves both left and right classes of morphisms componentwise for each $\Sigma$-functor symbol $B\colon \bar{c}\to d$. 
       \item If each $M_c$ is representably regular, and each arity support $M(t_1)$ and coarity support $M(t_2)$ representably preserve regular epimorphisms and finite limits, respectively, then $M^T$ is representably regular with $u$ representably preserving finite limits and regular epimorphisms. In particular, the conclusion follows if each $M_c$ is representably regular, $T$ is algebraic, and $M(B)$ representably preserves regular epis for each functor symbol $B$.
       \item If each $M_c$ is representably exact and $M(t_1)$ and $M(t_2)$ representably preserve coequalizers of kernel pairs and finite limits, respectively, for arity terms $t_1$ and coarity terms $t_2$, then $M^T$ is representably exact. In particular, if each $M_c$ is representably exact, $T$ is algebraic and $M(B)$ representably preserves reflexive coequalizers componentwise, then $M^T$ is exact.
   \end{enumerate}
\end{corollary}
\begin{proof}
   Follows from Corollaries $\ref{cor:pie_ofs_creation_2cat}$ and $\ref{cor:pie_exactness_creation_2cat}$. 
\end{proof}

Consider a symmetric monoidal category $A$ and denote by $\textbf{Comon}^{\mathrm{coc}}_A$ the category of internal cocommutative comonoids. The category of internal cocommutative Hopf algebras $\Hopf_A^{\mathrm{coc}}$ in $A$ is defined as the category $\textbf{Grp}(\textbf{Comon}^{\mathrm{coc}}_A)$ of internal groups within $\textbf{Comon}^{\mathrm{coc}}_A$. In the case where $A$ is the symmetric monoidal category of $k$-vector spaces or its dual, it has been shown that the category $\Hopf^{\mathrm{coc}}_A$ is exact (even semi-abelian) \cite{GRAN20194171, forsman2026semiabeliannessaffinegroupschemes}. The following corollary generalizes some exactness results found in \cite{bevilacqua2025coalgebraicmodelsomegagroups, gran2025hopfbracessemiabeliancategories} by showing that $(\textbf{Comon}^{\mathrm{coc}}_A)^T$ is exact for all algebraic theories $T$ extending the theory of groups, when $A$ is either the symmetric monoidal category of $k$-vector spaces or its dual for a field $k$:


\begin{corollary}
    Let $(\Sigma,\sigma)$ be a signature pair with algebraic $\sigma$-theories $T_1\subset T_2$. Assume $M$ is a $\Sigma$-model in $\CAT$, where the following holds for any category symbol $c$ and a functor symbol $B$:
    \begin{enumerate}
        \item $M_c$ is finitely complete.
        \item $M_c$ carries a fixed orthogonal factorization system which is preserved by each functor $M(B)$ componentwise.
        \item $M^{T_1}$ is a regular category whose regular epi-mono factorization system coincides with the induced orthogonal factorization system. 
    \end{enumerate}
    Then $M^{T_2}$ is regular. Furthermore, if $M^{T_1}$ is exact, each $M_c$ has reflexive coequalizers, and each $M(B)$ preserves reflexive coequalizers componentwise, then $M^{T_2}$ is exact. 
\end{corollary}
\begin{proof}
    Consider the forgetful functor $U\colon M^{T_2}\to M^{T_1}$, which is a conservative functor preserving finite limits from a finitely complete category while respecting the induced factorization systems. The regularity of $M^{T_2}$ is then clear. In the exact case, $M^{T_2}$ has reflexive coequalizers and $ U$ preserves them (Corollary~\ref{cor:liftingSifted}), and a similar argument to one in Corollary~\ref{cor:liftingregularexact} shows how the exactness is lifted from  $M^{T_1}$ to $M^{T_2}$. 
\end{proof}
\section*{Declarations}
\textbf{Funding} \\
\noindent
This research was supported by the Fonds de la Recherche Scientifique, FNRS (Belgium).

{\footnotesize
\begin{thebibliography}{10}

\bibitem{LPC}
J.~Adámek and J.~Rosický.
\newblock {\em Locally {Presentable} and {Accessible} {Categories}}, volume 189 of {\em London {Mathematical} {Society} {Lecture} {Note} {Series}}.
\newblock Cambridge University Press, Cambridge, UK, 1st edition, 1994.
\newblock URL: \url{https://www.cambridge.org/core/product/identifier/9780511600579/type/book}, \href {https://doi.org/10.1017/CBO9780511600579} {\path{doi:10.1017/CBO9780511600579}}.

\bibitem{AdamekRosickyVitale2010SiftedColimits}
J.~Adámek, J.~Rosický, and E.~M. Vitale.
\newblock What are sifted colimits?
\newblock {\em Theory and Applications of Categories}, 23(13):251--260, 2010.
\newblock URL: \url{http://www.tac.mta.ca/tac/volumes/23/13/23-13abs.html}.

\bibitem{agore2025categoryhopfbraces}
A.~Agore and A.~Chirvasitu.
\newblock On the category of {Hopf} braces, 2025.
\newblock arXiv: 2503.06280 [math.QA].
\newblock URL: \url{https://arxiv.org/abs/2503.06280}.

\bibitem{bevilacqua2025coalgebraicmodelsomegagroups}
M.~Bevilacqua.
\newblock Coalgebraic models of omega-groups, 2025.
\newblock arXiv: 2510.22653 [math.CT].
\newblock URL: \url{https://arxiv.org/abs/2510.22653}.

\bibitem{dubuc1970}
E.~J. Dubuc.
\newblock {\em Kan extensions in enriched category theory}, volume 145 of {\em Lecture {Notes} in {Mathematics}}.
\newblock Springer-Verlag, Berlin–Heidelberg–New York, 1970.
\newblock \href {https://doi.org/10.1007/BFb0060485} {\path{doi:10.1007/BFb0060485}}.

\bibitem{forsman2026semiabeliannessaffinegroupschemes}
D.~Forsman.
\newblock On the semi-abelianness of affine group schemes, 2026.
\newblock arXiv: 2602.21060 [math.CT].
\newblock URL: \url{https://arxiv.org/abs/2602.21060}.

\bibitem{GordonPowerStreet1995CoherenceTricategories}
R.~Gordon, A.~J. Power, and R.~Street.
\newblock {\em Coherence for tricategories}, volume 117 of {\em Memoirs of the {American} {Mathematical} {Society}}.
\newblock American Mathematical Society, Providence, RI, 1995.
\newblock Number: 558.
\newblock URL: \url{https://researchers.mq.edu.au/en/publications/coherence-for-tricategories/}, \href {https://doi.org/10.1090/memo/0558} {\path{doi:10.1090/memo/0558}}.

\bibitem{gran2025hopfbracessemiabeliancategories}
M.~Gran and A.~Sciandra.
\newblock Hopf braces and semi-abelian categories, 2025.
\newblock arXiv: 2411.19238 [math.RA].
\newblock URL: \url{https://arxiv.org/abs/2411.19238}.

\bibitem{GRAN20194171}
M.~Gran, F.~Sterck, and J.~Vercruysse.
\newblock A semi-abelian extension of a theorem by {Takeuchi}.
\newblock {\em Journal of Pure and Applied Algebra}, 223(10):4171--4190, 2019.
\newblock URL: \url{https://www.sciencedirect.com/science/article/pii/S002240491930012X}, \href {https://doi.org/10.1016/j.jpaa.2019.01.004} {\path{doi:10.1016/j.jpaa.2019.01.004}}.

\bibitem{GrandisTholen2006NaturalWFS}
M.~Grandis and W.~Tholen.
\newblock Natural weak factorization systems.
\newblock {\em Archivum Mathematicum}, 42(4):397--408, 2006.
\newblock URL: \url{http://eudml.org/doc/249802}.

\bibitem{Kelly:Enriched}
G.~M. Kelly.
\newblock {\em Basic concepts of enriched category theory}, volume~64 of {\em London {Mathematical} {Society} {Lecture} {Notes} {Series}}.
\newblock Cambridge University Press, 1982.
\newblock URL: \url{http://www.tac.mta.ca/tac/reprints/articles/10/tr10abs.html}.

\bibitem{Kelly1991}
G.~M. Kelly.
\newblock A note on relations relative to a factorization system.
\newblock In A.~Carboni, M.~C. Pedicchio, and G.~Rosolini, editors, {\em Category theory}, volume 1488 of {\em Lecture {Notes} in {Mathematics}}, pages 249--261. Springer, Berlin, Heidelberg, 1991.
\newblock URL: \url{https://link.springer.com/chapter/10.1007/BFb0084224}, \href {https://doi.org/10.1007/BFb0084224} {\path{doi:10.1007/BFb0084224}}.

\bibitem{LACK2023107196}
S.~Lack and G.~Tendas.
\newblock Virtual concepts in the theory of accessible categories.
\newblock {\em Journal of Pure and Applied Algebra}, 227(2):107196, 2023.
\newblock URL: \url{https://www.sciencedirect.com/science/article/pii/S002240492200192X}, \href {https://doi.org/10.1016/j.jpaa.2022.107196} {\path{doi:10.1016/j.jpaa.2022.107196}}.

\bibitem{Linton1969}
F.~E.~J. Linton.
\newblock Coequalizers in categories of algebras.
\newblock In B.~Eckmann, editor, {\em Seminar on triples and categorical homology theory}, volume~80 of {\em Lecture {Notes} in {Mathematics}}, pages 75--90. Springer, Berlin, Heidelberg, 1969.
\newblock URL: \url{https://link.springer.com/chapter/10.1007/BFb0083082}, \href {https://doi.org/10.1007/BFb0083082} {\path{doi:10.1007/BFb0083082}}.

\bibitem{MacLane1998}
S.~Mac~Lane.
\newblock {\em Categories for the working mathematician}, volume~5 of {\em Graduate {Texts} in {Mathematics}}.
\newblock Springer, Berlin, Heidelberg, 2nd edition, 1998.
\newblock URL: \url{https://link.springer.com/book/10.1007/978-1-4757-4721-8}, \href {https://doi.org/10.1007/978-1-4757-4721-8} {\path{doi:10.1007/978-1-4757-4721-8}}.

\bibitem{Porst01062008}
H.-E. Porst.
\newblock On categories of monoids, comonoids, and bimonoids.
\newblock {\em Quaestiones Mathematicae}, 31(2):127--139, 2008.
\newblock URL: \url{https://user.informatik.uni-bremen.de/porst/dvis/CoBiMon_Porst.pdf}, \href {https://doi.org/10.2989/QM.2008.31.2.2.474} {\path{doi:10.2989/QM.2008.31.2.2.474}}.

\bibitem{PorstHopfLocallyPresentable}
H.-E. Porst.
\newblock The formal theory of hopf algebras part {I}: {Hopf} monoids in a monoidal category.
\newblock {\em Quaestiones Mathematicae}, 38:1--52, June 2015.
\newblock URL: \url{https://user.informatik.uni-bremen.de/porst/dvis/PART%20I.pdf}, \href {https://doi.org/10.2989/16073606.2014.981736} {\path{doi:10.2989/16073606.2014.981736}}.

\bibitem{PowerPIE1991}
J.~Power and E.~Robinson.
\newblock A characterization of pie limits.
\newblock {\em Mathematical Proceedings of the Cambridge Philosophical Society}, 110:33 -- 47, July 1991.
\newblock \href {https://doi.org/10.1017/S0305004100070092} {\path{doi:10.1017/S0305004100070092}}.

\bibitem{Riehl2014CategoricalHomotopyTheory}
E.~Riehl.
\newblock {\em Categorical homotopy theory}, volume~24 of {\em New {Mathematical} {Monographs}}.
\newblock Cambridge University Press, Cambridge, 2014.
\newblock URL: \url{https://emilyriehl.github.io/files/cathtpy.pdf}, \href {https://doi.org/10.1017/CBO9781107261457} {\path{doi:10.1017/CBO9781107261457}}.

\bibitem{Ulmer1977Bialgebras}
F.~Ulmer.
\newblock Bialgebras in locally presentable categories.
\newblock Preprint, University of Wuppertal, 1977.
\newblock URL: \url{https://github.com/CategoryTheoryArchive/archive/blob/main/resources/1977_ulmer_bialgebras-locally.pdf}.

\bibitem{Wissmann2022Minimality}
T.~Wißmann.
\newblock Minimality notions via factorization systems and examples.
\newblock {\em Logical Methods in Computer Science}, 18(3):31:1--31:31, 2022.
\newblock URL: \url{https://lmcs.episciences.org/10014}, \href {https://doi.org/10.46298/lmcs-18(3:31)2022} {\path{doi:10.46298/lmcs-18(3:31)2022}}.

\end{thebibliography}
}
\end{document}